% !TEX program = pdflatex
\documentclass[12pt]{article}

\usepackage[utf8]{inputenc}
\usepackage{CJKutf8}
\usepackage{amsmath,amssymb}
\usepackage{amsthm}
\usepackage{geometry}
\geometry{
  a4paper,
  top=25mm,
  bottom=25mm,
  left=20mm,
  right=20mm
}
\usepackage{xcolor}
\usepackage{url}
\usepackage[unicode,colorlinks=true,linkcolor=blue,urlcolor=blue]{hyperref}
\usepackage{xurl}
\Urlmuskip=0mu plus 2mu

\setlength{\emergencystretch}{6em}

\newtheorem{definition}{定义}[section]
\newtheorem{theorem}{定理}[section]
\newtheorem{proposition}{命题}[section]
\newtheorem{remark}{备注}[section]

\newcommand{\JacGap}{\mathsf{JacGap}}
\newcommand{\GZLS}{\ensuremath{\mathsf{GZ\mbox{-}LS}}}

\begin{document}
\begin{CJK*}{UTF8}{gkai}
\sloppy

\title{雅可比间隙与 D 结构在语义层面的严格等价性（工作笔记）}
\author{Gao Zheng（高政）}
\date{2025-12-15}
\maketitle

\section*{前言}
\addcontentsline{toc}{section}{前言}

本文的目标不是单独比较 O3 理论中的 D 结构与 G-Framework 中的 Jacobiator-gap，
而是在语义层面给出二者之间的最小接口等价。若二者只在术语上并列，D 结构容易被
理解为决策偏好工具，JacGap 容易被理解为代数一致性残差；本文将它们统一到路径
偏好、逻辑性度量与证书型不变量的共同语义空间中。

本文的第一层立场是 D 结构的双重语义。D 既表示 Decision，即候选路径之间的偏好
与选择；也表示 Depth，即由自反性迭代更新诱导的层级深度。二者不是两套无关含义，
而是同一语义接口的两种投影：一个面向当前路径排序，一个面向更新进深和修复层级。

本文的第二层立场是 Jacobiator-gap 的证书语义。JacGap 不只是一个数值误差，而是
说明某个路径、组合或推理结构存在闭合缺口的证书型不变量。它能够诱导偏好关系：
间隙更小、闭合更好、证书更稳定的路径，在语义上优于间隙更大或证书不足的路径。

本文的第三层立场是双向构造。本文证明每个满足最小接口的 D 结构都可以嵌入到某种
JacGap 证书体系中；反过来，每个满足正则性的 JacGap 证书体系也可以实现为某个
D 结构所生成的逻辑性度量算子。二者在路径偏好关系这一统一语义空间中互为逆映射。

本文的第四层立场是路径偏好关系的统一。无论 D 结构从环境和历史中生成权重，还是
JacGap 从闭合缺口中生成证书，它们最终都在候选路径集合上产生可比较的偏好关系。
本文将这个偏好关系作为共同语义空间，避免在实现细节层面强行比较两套体系。

本文的第五层立场是逻辑性度量的接口化。逻辑性度量不是单纯的数值打分，而是把
环境、历史、维度投影和路径性质合成一个可执行排序规则。JacGap 体系若满足正则性，
也可以生成同样类型的排序规则，因此两者可以通过输出语义而非内部术语建立等价。

本文的第六层立场是证书型不变量的决策作用。证书不只是事后证明，它也可以参与
路径选择：一个路径的证书更强、缺口更小、闭合更稳定，就可以在决策语义中占优。
这使 JacGap 从“数学残差”进入“决策深度”的空间。

本文的第七层立场是为后续 G-Framework 主链提供早期决策接口。后续同伦扭曲、
修复塔、开放系统治理和高阶正交命中都需要比较路径或策略。本文给出的 D 结构与
JacGap 等价，正是这些比较过程的早期语义底座。

因此，本文在整体 G-Framework 中承担的是“决策深度与 JacGap 证书等价”的基础锚定。
它向前连接路径选择与逻辑性度量，向中间连接自反更新、修复深度和闭合缺口，向后
为同伦最小性、边缘算子、开放系统治理和高阶正交命中提供共同的决策语义底座。
概括地说，本文把“如何比较路径”推进为“D 结构与 JacGap 证书在语义层面严格等价”的
形式系统。

更具体地说，本文把路径选择从经验偏好提升为证书偏好。一个路径更好，不只是因为
它被某个模型打了更高分，而是因为在逻辑性度量或 JacGap 证书下，它表现出更强的
闭合性、更低的缺口或更稳定的更新深度。这样，决策与修复之间获得了共同语言。

从方法论上看，本文是 G-Framework 中“路径如何排序”的早期答案。后续同伦扭曲、
策略迁移、开放系统治理和统计命中都需要在候选路径之间做选择；本文说明这种选择
可以被压缩为 D 结构和 JacGap 证书之间的语义等价。

从整体谱系看，本文把外部理论接口与 G-Framework 内部证书体系接合起来。它既保留
O3 D 结构的决策/深度含义，又把它翻译成 JacGap 证书语言，使后续体系可以在不丢失
决策语义的前提下进入同伦修复和高阶控制框架。

\tableofcontents
\clearpage

\section{写作目的与背景说明}

本笔记的目标，是在 \emph{语义层面}（而不是仅仅在符号或实现层面）
精确定义并证明：
\begin{center}
  \textbf{O3 理论中的 D 结构，与 G-Framework 中由雅可比间隙（Jacobiator-gap）
  驱动的证书型不变量，其“决策语义”是严格等价的。}
\end{center}

此外，本文对 “D” 的命名采取\emph{双重语义}对齐：一方面它对应
\textbf{Decision}（对路径的偏好/选择语义），另一方面它也对应
\textbf{Depth}（由自反性迭代更新所诱导的层级/序数高度）。
前者用于给出可执行的路径偏好关系，后者用于刻画“更新进深”（例如后文的
超限递归修复深度）。两者将在第 \ref{sec:D-dual} 节中形式化并证明为
同一 D 结构语义接口的两种投影。

更具体地，我们希望证明如下命题：
\begin{quotation}
  在合适的抽象下，\textbf{每一个由 D 结构生成的逻辑性度量算子
  所定义的路径偏好关系，都可以唯一嵌入到一种雅可比间隙证书体系中；}
  反之，\textbf{每一个满足一定正则性的 Jacobiator-gap 证书体系，
  都可以被实现为某个 D 结构所生成的逻辑性度量算子。}
\end{quotation}

我们所使用的两个主要参考源为：
\begin{itemize}
  \item O3 理论：D 结构与逻辑性度量（kernel\_reference）\cite{O3DStructureLogicalMetric}；
  \item G-Framework 纯数学发布稿 \cite{GaoZheng2025GFramework}（引入 law-space、三层联络与 Jacobiator-gap 证书型不变量 $\JacGap$）。
\end{itemize}

本笔记的工作方式是：
\begin{enumerate}
  \item 抽象出 D 结构及其逻辑性度量的\emph{最小语义接口}；
  \item 抽象出 Jacobiator-gap 证书体系的\emph{最小语义接口}；
  \item 给出一个双向的构造：从 D 结构到 Jacobiator-gap，从 Jacobiator-gap
        到 D 结构；
  \item 在“路径偏好关系”这一统一语义空间中，证明这两个构造互为逆映射
        （在自然的等价关系下），从而得到“语义上的严格等价性”。
\end{enumerate}

\section{D 结构与逻辑性度量的语义抽象}

\subsection{D 结构的最小语义接口}

根据 \cite{O3DStructureLogicalMetric}，
我们可以把 D 结构抽象为如下数据：
\begin{itemize}
  \item 一个环境/历史空间 $(E,H)$，其中 $E$ 表示当前客观环境
        与外部“压强吸引子”，$H$ 表示系统的经验历史；
  \item 一个由 DERI 算法给出的\emph{基准拟合器}
        \[
          \mathsf{Fit}_D : (E,H)\longrightarrow w,
        \]
        输出一个权重向量 $w=(w_1,\dots,w_n)$，刻画当前“价值观基准”；
  \item 一个逻辑性度量算子
        \[
          \mathcal{L}_D(\gamma;w)
          =
          \sum_{i=1}^n w_i\cdot d_i(\gamma),
        \]
        其中 $\gamma$ 是一条候选路径（决策、控制或演化轨迹），
        $d_i(\gamma)$ 是在维度 $i$ 上对路径的“逻辑性”投影。
\end{itemize}
在这个抽象下，D 结构的\textbf{输出语义}可以简化为：
\begin{equation}
  \gamma_1
  \succ_D
  \gamma_2
  \quad\Longleftrightarrow\quad
  \mathcal{L}_D(\gamma_1;w)
  >
  \mathcal{L}_D(\gamma_2;w),
  \label{eq:D-order}
\end{equation}
即 D 结构通过 $\mathcal{L}_D$ 在候选路径集合上定义了一个
严格偏好关系 $\succ_D$。

\subsection{自反性与迭代更新}

同一文档中强调，D 结构具有\emph{自反性（reflexivity）}：
系统执行路径 $\gamma$ 后，会产生新的经验并更新 $(E,H)$，
进而通过 $\mathsf{Fit}_D$ 生成新的基准 $w'$。这一点在语义上
可以概括为：
\[
  (E,H) \xrightarrow{\ \mathsf{Fit}_D\ } w
  \xrightarrow{\ \gamma\ } (E',H') \xrightarrow{\ \mathsf{Fit}_D\ } w'.
\]

在本笔记中，我们主要关心的是：\emph{在某一时刻固定的 $w$ 下，
由式 \eqref{eq:D-order} 定义的偏好关系}。自反性保证了这一偏好关系
可以随时间演化，但在每一个“冻结”时刻，我们都可以视之为
一个静态的决策语义。

\subsection{D 的双重语义：Depth 与 Decision}\label{sec:D-dual}

本节把 D 结构的“决策”与“深度”两层含义都显式写为可检验的数学对象。
为避免把命名当作修辞，我们只做两件事：
\begin{itemize}
  \item 把 \eqref{eq:D-order} 给出的路径偏好关系视为 \textbf{D-Decision} 语义；
  \item 把自反性更新的迭代层级（必要时提升为序数）视为 \textbf{D-Depth} 语义。
\end{itemize}

\begin{definition}[D-Decision：路径偏好与最大元集合]
给定固定的 D 结构状态 $(E,H,w)$，令 $\Gamma(E,H)$ 表示在该状态下的可行路径集合。
由 \eqref{eq:D-order} 得到严格偏好关系 $\succ_D$。
定义（可能为集合值的）决策输出为
\[
  \mathrm{Dec}_D(E,H,w)
  :=
  \bigl\{
    \gamma\in\Gamma(E,H)\ \big|\ \nexists\gamma'\in\Gamma(E,H)\ \text{s.t.}\ \gamma'\succ_D\gamma
  \bigr\}.
\]
这一定义不要求“唯一最优”，只要求能在偏好关系下识别最大元集合。
\end{definition}

\begin{definition}[D-Depth：更新塔与序数高度]
把自反性更新抽象为一个更新算子
\[
  \mathsf{Upd}_D:\ (E,H,w)\times\Gamma(E,H)\longrightarrow (E',H',w').
\]
给定初始状态 $D_0=(E_0,H_0,w_0)$，取任意选择规则使得
$\gamma_n\in\mathrm{Dec}_D(E_n,H_n,w_n)$，并递推定义
\[
  D_{n+1}:=\mathsf{Upd}_D(D_n,\gamma_n).
\]
若存在一个取值于良序集合（例如序数）的\emph{排名函数}
$\rho$，满足对所有允许的更新都有严格下降
\[
  \rho(D_{n+1})<\rho(D_n),
\]
则称该迭代过程具有良基“深度”，并把其\emph{序数高度}
（最小无法继续下降的临界层级）记为 $\mathrm{Depth}_D(D_0)$。
\end{definition}

\begin{theorem}[D 结构的双重语义性：Decision + Depth]\label{thm:D-dual-semantics}
在上述抽象下，任意一个带自反性更新的 D 结构同时携带：
\begin{enumerate}
  \item \textbf{决策语义（Decision）：}由 $\mathcal{L}_D$ 诱导的偏好关系 $\succ_D$（及其最大元集合 $\mathrm{Dec}_D$）；
  \item \textbf{深度语义（Depth）：}由更新算子 $\mathsf{Upd}_D$ 的迭代层级（在存在良基排名函数 $\rho$ 时可提升为确定的序数高度 $\mathrm{Depth}_D$）。
\end{enumerate}
因此，“D” 可以同时被理解为 \textnormal{Decision} 与 \textnormal{Depth} 的统一记号：它们是同一套 D 结构数据在
“路径排序/选择”与“迭代层级/进深”两个坐标上的投影。
\end{theorem}

\begin{proof}
（1）决策语义直接来自第 \ref{eq:D-order} 式：给定 $(E,H,w)$，
$\mathcal{L}_D(\cdot;w)$ 诱导严格偏好关系 $\succ_D$，从而确定
最大元集合 $\mathrm{Dec}_D(E,H,w)$。

（2）深度语义来自自反性：一旦给出更新算子 $\mathsf{Upd}_D$，
就得到一个按步数 $n$ 索引的更新塔 $\{D_n\}_{n\in\mathbb{N}}$；
若进一步存在良基排名函数 $\rho$ 使得每一步更新严格下降，
则不存在无限下降链，从而该塔在良序意义下具有确定的“高度”，记为
$\mathrm{Depth}_D(D_0)$。

以上两点表明：同一 D 结构数据既决定了对路径的排序/选择（Decision），
也决定了更新进深的层级刻度（Depth），因此具备双重语义。
\end{proof}

\begin{remark}[与后文超限模型的对齐]
后文在讨论“超限递归 + 同伦扭曲”时，会把更新塔的层级从自然数提升到序数，
记作 $\{D_\alpha\}_{\alpha<\kappa}$ 并以临界序数 $\alpha_\ast$ 表示“修复深度”。
这正是上面 $\mathrm{Depth}_D$ 的一个具体化：把排名函数 $\rho$ 取为
“剩余修复阶”的序数高度即可。
\end{remark}

\subsubsection*{补充：Depth 的本体论层级与递归序数刻度}

\begin{remark}[Depth 的含义说明]
本节的“深度（Depth）”指
\emph{本体论深度（ontological depth）}：
它度量系统在递归更新中\emph{需要上升到多高的语义层级}，
以便把“现象层的对象”提升为“规律层/元规律层的对象”。
这与机器学习语境中“网络层数（layer depth）”是不同概念。
\end{remark}

\begin{definition}[本体论层级：$L_0/L_1/L_\omega$]\label{def:onto-level}
用一族按序数索引的层级标签 $\{\mathsf{L}_\alpha\}_{\alpha\le\omega}$ 表示
“描述与解释的本体论层级”，其中
\[
  \mathsf{L}_0:\ \text{现象层（例如价格序列、生物表型等可直接观测量）},
\]
\[
  \mathsf{L}_1:\ \text{规律层（例如物理定律、经济模型等可压缩的生成规则）},
\]
\[
  \mathsf{L}_\omega:\ \text{超限层/元规则层（例如用于约束规律本身的法空间几何/证书结构）}.
\]
把“从 $\mathsf{L}_\alpha$ 上升到 $\mathsf{L}_{\alpha+1}$”理解为一次
\emph{解释提升（lifting）}：增加可用的结构词汇（不变量、证书、同伦约束等），以消解在较低层级无法闭合的障碍。
\end{definition}

\begin{definition}[超限递归更新与不动点]\label{def:transfinite-fixed-point}
把“$D\to D'$ 的递归深化”抽象为一个更新算子
$\mathsf{Upd}_D$（见前文定义），并在需要时把迭代提升为序数索引：
给定 $D_0$，对任意序数 $\alpha$ 递推定义
\[
  D_{\alpha+1}:=\mathsf{Upd}_D(D_\alpha,\gamma_\alpha),
\qquad
  \gamma_\alpha\in \mathrm{Dec}_D(D_\alpha),
\]
并在极限序数 $\lambda$ 处用某种汇总规则（例如并/上确界/余极限）定义 $D_\lambda$。
给定一个语义等价关系 $\simeq$（表示“在路径偏好意义下等价”），称 $D_\alpha$ 为\emph{不动点}，若
\[
  D_{\alpha+1}\ \simeq\ D_\alpha.
\]
\end{definition}

\begin{proposition}[Depth 等价于最小不动点序数]\label{prop:depth-fixed-point}
在定义 \ref{def:transfinite-fixed-point} 的设定下，若存在最小序数 $\alpha_\ast$
使得 $D_{\alpha_\ast}$ 为不动点，则可把
\[
  \mathrm{Depth}_D(D_0):=\alpha_\ast
\]
作为“本体论深度”的一种等价刻度；它与前文基于良基排名函数 $\rho$ 的
深度定义一致（在把 $\rho$ 取为“距离不动点的剩余阶”时）。
\end{proposition}

\begin{proof}
若存在最小不动点序数 $\alpha_\ast$，则序列在 $\alpha_\ast$ 处稳定：
$D_{\alpha_\ast+1}\simeq D_{\alpha_\ast}$。
此时可定义排名函数 $\rho$ 为
\[
  \rho(D_\alpha):=\min\{\beta\mid D_{\alpha+\beta}\simeq D_{\alpha_\ast}\},
\]
它取值于序数并在每一步更新中严格下降，因而满足前文“良基深度”的要求，
并且其临界高度正是 $\alpha_\ast$。反过来，若已给出一个“剩余修复阶”式的
良基排名函数 $\rho$，则当 $\rho(D_\alpha)=0$ 时系统在语义上无法再下降，
等价于达到不动点。两种刻度互相重标度，因此一致。
\end{proof}

\begin{remark}[递归序数刻度（可执行版本）]
若进一步要求更新过程是\emph{可执行的}（例如每一步由 DERI 等算法产生），则
我们只能沿着可计算的序数记号系统推进：也就是说，允许的层级索引
$\alpha$ 必须能被一个递归记号表示（常称为\emph{递归序数}）。
在该可执行语义下，“Depth $\leftrightarrow$ 递归序数”并非额外结论，
而是对“可执行的超限递归”所能触及的层级边界的直接刻画。
\end{remark}

\subsubsection*{补充：Decision 作为选择算子、剪枝与切割}

\begin{definition}[选择函数（choice operator）]\label{def:choice-operator}
给定集合 $\Gamma$，称映射
\[
  \mathsf{Ch}:\ \mathcal{P}(\Gamma)\setminus\{\varnothing\}\longrightarrow \Gamma
\]
为一个选择函数，若对任意非空子集 $A\subseteq \Gamma$ 都有
$\mathsf{Ch}(A)\in A$。
\end{definition}

\begin{proposition}[Decision 输出到“单一路径”的提升]\label{prop:decision-choice}
在 D-Decision 定义下，$\mathrm{Dec}_D(E,H,w)$ 可能是集合值。
若给定一个选择函数 $\mathsf{Ch}$，则可定义单值决策算子
\[
  \mathrm{Dec}_D^{\mathsf{Ch}}(E,H,w):=\mathsf{Ch}\bigl(\mathrm{Dec}_D(E,H,w)\bigr),
\]
从而把“偏好关系”提升为“可执行的单一路径选择”。
\end{proposition}

\begin{proof}
由定义 \ref{def:choice-operator}，对任意非空集合 $A$ 都有 $\mathsf{Ch}(A)\in A$。
当 $\mathrm{Dec}_D(E,H,w)\ne\varnothing$ 时，取 $A=\mathrm{Dec}_D(E,H,w)$，
即得 $\mathrm{Dec}_D^{\mathsf{Ch}}(E,H,w)\in \mathrm{Dec}_D(E,H,w)$，从而是一个合法的最大元选择。
\end{proof}

\begin{theorem}[选择公理的路径空间版本（等价式）]\label{thm:ac-path}
对固定的路径集合 $\Gamma$，下列两条等价：
\begin{enumerate}
  \item （路径空间上的选择公理）对任意指标集 $I$ 及任意一族非空子集 $\{A_i\subseteq\Gamma\}_{i\in I}$，
        存在函数 $f:I\to\Gamma$ 使得 $f(i)\in A_i$；
  \item （全域选择函数）存在选择函数 $\mathsf{Ch}:\mathcal{P}(\Gamma)\setminus\{\varnothing\}\to\Gamma$。
\end{enumerate}
\end{theorem}

\begin{proof}
（2）$\Rightarrow$（1）：令 $f(i):=\mathsf{Ch}(A_i)$ 即可。

（1）$\Rightarrow$（2）：取指标集 $I:=\mathcal{P}(\Gamma)\setminus\{\varnothing\}$，
并对每个 $A\in I$ 令 $A_A:=A$。由（1）存在 $f:I\to\Gamma$ 使得 $f(A)\in A$。
令 $\mathsf{Ch}(A):=f(A)$，即得（2）。
\end{proof}

\begin{definition}[剪枝切割（Dedekind cut 的评分版本）]\label{def:dedekind-cut}
给定评分函数 $f:\Gamma\to\mathbb{R}$，它诱导一个预序
\[
  \gamma_1\preceq_f \gamma_2 \quad:\Longleftrightarrow\quad f(\gamma_1)\le f(\gamma_2).
\]
对任意阈值 $\tau\in\mathbb{R}$，定义切割（剪枝）：
\[
  \Gamma_{<\tau}:=\{\gamma\in\Gamma\mid f(\gamma)<\tau\},
\qquad
  \Gamma_{\ge\tau}:=\{\gamma\in\Gamma\mid f(\gamma)\ge\tau\}.
\]
\end{definition}

\begin{proposition}[切割诱导的“信号/噪声”划分是模型内可判定谓词]\label{prop:signal-noise-cut}
在定义 \ref{def:dedekind-cut} 下，二元谓词
\[
  \mathrm{Signal}_\tau(\gamma)\iff \gamma\in\Gamma_{\ge\tau},
\qquad
  \mathrm{Noise}_\tau(\gamma)\iff \gamma\in\Gamma_{<\tau}
\]
给出一个由评分函数 $f$ 完全决定的划分。
因此，“剪枝”可以被严格写成：把候选集合替换为 $\Gamma_{\ge\tau}$，再在其上做选择（见命题 \ref{prop:decision-choice}）。
\end{proposition}

\begin{proof}
由定义 \ref{def:dedekind-cut}，对任意 $\gamma\in\Gamma$，要么 $f(\gamma)<\tau$，
要么 $f(\gamma)\ge\tau$，从而 $\Gamma=\Gamma_{<\tau}\,\sqcup\,\Gamma_{\ge\tau}$。
谓词由 $f$ 与 $\tau$ 唯一确定。将候选集合替换为 $\Gamma_{\ge\tau}$ 即为剪枝；随后用选择函数从剩余集合中取代表元即可。
\end{proof}

\begin{remark}[“选择哪一套逻辑有效”的一种形式化方式]
在更高一层的建模中，可以把 $\Gamma$ 取为“可选的公理/规则配置”的集合，
并把评分函数 $f$ 解释为“在给定战略目标与一致性约束下的可行性/效用”。
此时，剪枝切割对应于先排除不满足一致性或信息增益阈值的配置（噪声），
再在剩余配置中执行选择（信号）。该形式化强调：Decision 不仅是“挑一个最优解”，
也可以是“先裁定可用的逻辑扇区，再在扇区内择优”。
\end{remark}

\section{Jacobiator-gap 证书体系的语义抽象}

\subsection{law-space 与三层联络的最小背景}

在 \cite{GaoZheng2025GFramework} 中，G-Framework
通过 law-space $\mathcal{L}$ 与三层联络（GZ-TLC）来刻画法则空间中的
几何与同伦结构。对本笔记而言，我们只需要以下简化数据：
\begin{itemize}
  \item 一个法则空间 $\mathcal{L}$，其点 $L\in\mathcal{L}$ 表示
        “系统当前所采用的法则组合”；
  \item 一个三层系统 $S$，它在 law-space 中决定了允许的演化轨迹
        与 curvature 结构；
  \item 对于每个 $S$，存在一个与之关联的“law-level 系统”
        $L(S)$，其代数结构在一般情况下只满足\emph{同伦 Jacobi 性}。
\end{itemize}

\subsection{Jacobiator-gap 作为证书型不变量}

在该主稿中，引入了一个典型的证书型不变量——雅可比间隙：
\[
  \JacGap(L)
  =
  \bigl(c_1(L),\dots,c_k(L)\bigr),
\]
大致满足如下性质：
\begin{itemize}
  \item $c_i(L)$ 度量了在不同层面上对严格 Jacobi 恒等式的
        控制性违反程度；
  \item $\JacGap(L)=0$ 当且仅当 $L$ 可以在不改变低阶观测量的前提下，
        被提升为严格 Jacobi 类型的 law-object；
  \item $\JacGap$ 对 law-space 中由三层联络诱导的变形是稳定的：
        在合适意义下，沿着受控的同伦变形，$\JacGap$ 连续变化。
\end{itemize}

当一个三层系统 $S$ 通过 $L(S)$ 将自身投影到 law-level 后，我们得到
一个组合证书
\[
  S \longmapsto \JacGap\bigl(L(S)\bigr),
\]
这就是主稿中所说的 Jacobiator-gap 证书体系。

为了与 D 结构进行比较，我们从中抽象出如下\emph{语义接口}：
\begin{equation}
  \gamma_1
  \succ_{\JacGap}
  \gamma_2
  \quad\Longleftrightarrow\quad
  F\bigl(\JacGap(L),\gamma_1\bigr)
  >
  F\bigl(\JacGap(L),\gamma_2\bigr),
  \label{eq:Jac-order}
\end{equation}
其中 $F$ 是某个由主稿中的 law-curvature 与 Jacobiator-gap 推导出的
“路径证书函数”，它对每条路径 $\gamma$ 给出一个可审计的数值评分。

换言之，Jacobiator-gap 证书体系在语义上同样通过一个“分数函数”
诱导了路径上的偏好关系 $\succ_{\JacGap}$。

\section{统一的“语义等价”定义}

\subsection{路径偏好关系作为统一语义载体}

为了比较两套看似完全不同的数学结构，我们需要一个统一的语义载体。
在决策与控制场景中，一个最自然的选择是：\emph{路径偏好关系}。

设 $\Gamma$ 是某一时刻系统所有可行路径的集合。对任意两条路径
$\gamma_1,\gamma_2\in\Gamma$，我们关心的不是它们的具体数值评分，
而是：
\[
  \gamma_1 \text{ 是否被认为“更优于” } \gamma_2.
\]

形式化地说，一个\textbf{路径语义}可以被表示为一个严格偏好关系
$\succ$，满足：
\begin{itemize}
  \item 反自反性：$\gamma\not\succ\gamma$；
  \item 传递性：若 $\gamma_1\succ\gamma_2$ 且 $\gamma_2\succ\gamma_3$，
        则 $\gamma_1\succ\gamma_3$；
  \item 与某个实值评分函数 $f:\Gamma\to\mathbb{R}$ 相容，即
        $\gamma_1\succ\gamma_2 \Leftrightarrow f(\gamma_1)>f(\gamma_2)$。
\end{itemize}

在这种设定下，我们可以说：
\begin{itemize}
  \item D 结构给出了一个偏好关系 $\succ_D$（由 $\mathcal{L}_D$ 诱导）；
  \item Jacobiator-gap 证书体系给出了一个偏好关系 $\succ_{\JacGap}$
        （由 $F\circ\JacGap$ 诱导）。
\end{itemize}

\subsection{基本假设与公理化条件}

为了在后续论证中避免隐含使用条件，我们在此显式列出若干
关于路径语义与评分函数的基本假设：
\begin{itemize}
  \item[(A1)] \textbf{评分表示假设：}
    对于每一个路径偏好关系 $\succ$，存在实值函数
    $f:\Gamma\to\mathbb{R}$，使得
    $\gamma_1\succ\gamma_2 \Leftrightarrow f(\gamma_1)>f(\gamma_2)$。
  \item[(A2)] \textbf{单调重标度假设：}
    若两个评分函数 $f_1,f_2:\Gamma\to\mathbb{R}$ 诱导同一偏好关系
    $\succ$，则存在严格单调递增函数 $\phi:\mathbb{R}\to\mathbb{R}$，
    使得 $f_2=\phi\circ f_1$。
  \item[(A3)] \textbf{同一语义空间假设：}
    在比较 D 结构与 Jacobiator-gap 证书体系时，它们作用在同一条
    路径集合 $\Gamma$ 上，且所有评分函数都以 $\Gamma$ 为定义域。
\end{itemize}

在 (A1)–(A3) 的前提下，后文对“语义等价”的定义与双向构造定理
可以在统一的路径语义空间内被精确表述。

\subsection{语义等价性的定义}

\begin{definition}[路径偏好意义下的语义等价]
给定两个路径语义
  $\succ^{(1)}$ 与 $\succ^{(2)}$，若存在一个\emph{严格单调递增}的函数
  $\phi:\mathbb{R}\to\mathbb{R}$，以及两个评分函数
  $f_1,f_2:\Gamma\to\mathbb{R}$，使得
  \[
    \forall\gamma_1,\gamma_2\in\Gamma,\quad
    \gamma_1\succ^{(1)}\gamma_2
    \;\Longleftrightarrow\;
    f_1(\gamma_1)>f_1(\gamma_2),
  \]
  且
  \[
    f_2(\gamma)
    =\phi\bigl(f_1(\gamma)\bigr)
    \quad\forall\gamma\in\Gamma,
  \]
  则称 $\succ^{(1)}$ 与 $\succ^{(2)}$ 在\emph{路径偏好意义下语义等价}。
\end{definition}

直观而言，这表示两套语义系统在所有可选路径上给出的排序完全一致，
差别仅在于“打分的刻度不同”，而刻度之间是严格单调的可逆变换。

记 D 结构与 Jacobiator-gap 所给出的偏好关系分别为
$\succ_D$ 与 $\succ_{\JacGap}$。我们的目标就是证明：
\[
  \succ_D \quad\text{与}\quad \succ_{\JacGap}
  \quad\text{在上述意义下语义等价。}
\]

\subsection{路径语义范畴的公理化}

基于假设 (A1)–(A3)，可以把“路径 + 偏好 + 评分函数”的数据进一步
封装成一个范畴化的语义对象，从而为后文讨论“函子级别的等价”预留接口。

\begin{definition}[路径语义对象]
一个\emph{路径语义对象} $X$ 定义为三元组
  \[
    X=\bigl(\Gamma_X,\succ_X,[f_X]\bigr),
  \]
  其中：
  \begin{itemize}
    \item $\Gamma_X$ 是路径集合；
    \item $\succ_X$ 是 $\Gamma_X$ 上的严格偏好关系；
    \item $[f_X]$ 是所有与 $\succ_X$ 相容的评分函数
          $f:\Gamma_X\to\mathbb{R}$ 的等价类，等价关系由 (A2) 给出：
          若存在严格单调递增的 $\phi$ 使得 $f_2=\phi\circ f_1$，
          则 $f_1\sim f_2$。
  \end{itemize}
  假设 (A1) 保证每个 $(\Gamma_X,\succ_X)$ 至少有一个代表元
  $f_X\in[f_X]$，假设 (A2) 保证该等价类是良定义的。
\end{definition}

\begin{definition}[路径语义范畴 $\mathsf{PathSem}$]
定义范畴 $\mathsf{PathSem}$ 如下：
  \begin{itemize}
    \item 对象为所有路径语义对象 $X=(\Gamma_X,\succ_X,[f_X])$；
    \item 对于两个对象 $X,Y$，一个\emph{路径语义态射}
          $T:X\to Y$ 是一个映射
          $T:\Gamma_X\to\Gamma_Y$，满足：
          \begin{itemize}
            \item[(PS1)] \textbf{偏好保持性：}
              若 $\gamma_1\succ_X\gamma_2$，则
              $T(\gamma_1)\succ_Y T(\gamma_2)$；
            \item[(PS2)] \textbf{评分相容性：}
              存在严格单调递增函数 $\phi_T:\mathbb{R}\to\mathbb{R}$，
              以及代表元 $f_X\in[f_X],f_Y\in[f_Y]$，使得
              \[
                f_Y\bigl(T(\gamma)\bigr)
                =
                \phi_T\bigl(f_X(\gamma)\bigr)
                \quad\forall\gamma\in\Gamma_X.
              \]
          \end{itemize}
    \item 恒等态射由恒等映射 $\mathrm{id}_{\Gamma_X}$ 给出，
          其对应的 $\phi_{\mathrm{id}}$ 取为恒等函数；
    \item 若 $T:X\to Y$、$S:Y\to Z$ 为态射，则其复合
          $S\circ T:X\to Z$ 定义为映射复合
          $\Gamma_X\xrightarrow{T}\Gamma_Y\xrightarrow{S}\Gamma_Z$。
  \end{itemize}
\end{definition}

易见 (PS1) 保证了态射在偏好层面上是单调的，而 (PS2) 则在评分函数层面
实现了“严格单调重标度”的相容性；由于严格单调函数在复合下封闭，
$\mathsf{PathSem}$ 的态射在复合下同样满足 (PS1)–(PS2)，从而确实构成
一个范畴。

在这一记号下，可以把本节“路径偏好意义下的语义等价”重新表述为：
给定两个路径语义对象 $X,Y$，若存在一个\emph{路径语义同构}
（即在 $\mathsf{PathSem}$ 中的可逆态射）$T:X\to Y$，
则称 $X$ 与 $Y$ 在路径语义意义下等价。
前文给出的关于 D 结构与 Jacobiator-gap 的等价性定理，
在引入适当的“系统范畴”后，可以进一步提升为
函子 $F_D,F_{\JacGap}:\mathsf{Sys}\to\mathsf{PathSem}$
之间的自然同构，这将留待未来更高层次的范畴化工作中展开。

\subsection{系统范畴 $\mathsf{Sys}$ 的最小构型}

在 $\mathsf{PathSem}$ 已经给出之后，我们可以把“携带 D 结构 /
law-level 表述的系统”抽象为一个范畴 $\mathsf{Sys}$，其对象是
带有两种语义呈现但在路径语义上等价的系统，态射则是同时保留
两种语义的映射。

\begin{definition}[系统对象]
一个\emph{系统对象} $S\in\mathsf{Sys}$ 由如下数据构成：
  \[
    S
    =
    \bigl(
      \Gamma_S,\;
      D_S,\;
      L_S
    \bigr),
  \]
  其中：
  \begin{itemize}
    \item $\Gamma_S$ 是该系统下可行路径的集合；
    \item $D_S$ 是定义在 $\Gamma_S$ 上的一套 D 结构数据
          （如 $(E_S,H_S,w_S,\mathcal{L}_S)$），从而按照前文规则
          诱导出一个路径语义对象
          \[
            F_D(S)
            =
            \bigl(\Gamma_S,\succ^{(D)}_S,h^{(D)}_S\bigr)
            \in\mathsf{PathSem};
          \]
    \item $L_S$ 是定义在 law-space 上的 law-level 系统数据\\
          （如 $L_S\in\mathcal{L}$ 及其 Jacobiator-gap 证书
          $\JacGap(L_S)$）。
          由此可以定义另一个路径语义对象 $F_{\JacGap}(S)$，令
          {\small
          \[
            F_{\JacGap}(S)
            =
            \bigl(\Gamma_S,\succ^{(\JacGap)}_S,h^{(\JacGap)}_S\bigr)
            \in\mathsf{PathSem}.
          \]
          }
  \end{itemize}
  其中 $h^{(D)}_S$ 与 $h^{(\JacGap)}_S$ 分别代表评分函数等价类
  $[f^{(D)}_S]$ 与 $[f^{(\JacGap)}_S]$。
  此外，我们要求存在一个在 $\mathsf{PathSem}$ 中的同构
  \[
    \iota_S:
    F_D(S)
    \xrightarrow{\;\cong\;}
    F_{\JacGap}(S),
  \]
  且其底层映射是 $\Gamma_S$ 上的恒等映射。换言之，$D_S$ 与 $L_S$
  必须在同一条路径集合上给出语义等价、仅相差严格单调重标度的
  路径偏好关系。
\end{definition}

\begin{definition}[系统态射]
给定两个系统对象
  $S=(\Gamma_S,D_S,L_S)$
  与
  $S'=(\Gamma_{S'},D_{S'},L_{S'})$，
  一个\emph{系统态射} $\Phi:S\to S'$ 定义为一个映射
  \[
    T_\Phi:\Gamma_S\longrightarrow\Gamma_{S'},
  \]
  满足：
  \begin{itemize}
    \item[(SM1)] \textbf{D 语义保留：}
      存在 $\mathsf{PathSem}$ 中的态射
      \[
        F_D(\Phi):
        F_D(S)\longrightarrow F_D(S'),
      \]
      其底层映射为 $T_\Phi$，并满足 (PS1)–(PS2)；
    \item[(SM2)] \textbf{Jacobiator-gap 语义保留：}
      存在 $\mathsf{PathSem}$ 中的态射
      \[
        F_{\JacGap}(\Phi):
        F_{\JacGap}(S)\longrightarrow F_{\JacGap}(S'),
      \]
      其底层映射同样为 $T_\Phi$，并满足 (PS1)–(PS2)；
    \item[(SM3)] \textbf{语义兼容性：}
      两种语义下的态射与对象上的同构相容，即在 $\mathsf{PathSem}$
      中满足
      \[
        F_{\JacGap}(\Phi)\circ\iota_S
        =
        \iota_{S'}\circ F_D(\Phi).
      \]
  \end{itemize}
\end{definition}

在上述定义下，把 $F_D$ 与 $F_{\JacGap}$ 看作从 $\mathsf{Sys}$ 到
$\mathsf{PathSem}$ 的函子，是自然的：
\begin{itemize}
  \item 在对象上，$F_D,F_{\JacGap}$ 分别取系统的 D 语义与 Jacobiator-gap
        语义；
  \item 在态射上，$F_D(\Phi),F_{\JacGap}(\Phi)$ 由 (SM1)、(SM2) 保证存在，
        且在复合与恒等下保持范畴结构。
\end{itemize}

进一步地，条件 (SM3) 表明，对每个系统对象 $S$，同构
$\iota_S:F_D(S)\to F_{\JacGap}(S)$ 在 $\mathsf{Sys}$ 的态射下是自然的。
换言之，$\{\iota_S\}_S$ 给出了两个函子
$F_D,F_{\JacGap}:\mathsf{Sys}\to\mathsf{PathSem}$ 之间的一个
自然同构，这就是前文“雅可比间隙与 D 结构语义等价性”定理在
范畴层面的抽象版本。

\subsubsection*{一个具体的系统对象与系统态射示例}

为了把上述抽象定义具体化，我们构造一个极简的系统对象
$S\in\mathsf{Sys}$，以及一个表示“同伦扭曲前后状态对比”的
系统态射 $\Phi:S\to S'$。

\paragraph{系统对象 $S$ 的构造}

设可行路径集合为一个两点集合
\[
  \Gamma_S=\{\gamma_{\mathrm{safe}},\gamma_{\mathrm{risky}}\},
\]
其中 $\gamma_{\mathrm{safe}}$ 表示“保守路径”，
$\gamma_{\mathrm{risky}}$ 表示“激进路径”。

\begin{itemize}
  \item \textbf{D 结构部分 $D_S$：}
        取固定环境/历史 $(E_S,H_S)$ 与权重向量 $w_S$，
        在本示例中不显式展开，只关心其诱导的逻辑性度量
        $\mathcal{L}_S$。定义评分函数
        \[
          f^{(D)}_S(\gamma_{\mathrm{safe}})=0,
          \qquad
          f^{(D)}_S(\gamma_{\mathrm{risky}})=1.
        \]
        根据 (A1)，这唯一确定了一个偏好关系
        $\succ^{(D)}_S$，即
        $\gamma_{\mathrm{risky}}\succ^{(D)}_S\gamma_{\mathrm{safe}}$。
        相应的路径语义对象为
        \[
          F_D(S)
          =
          \bigl(
            \Gamma_S,\succ^{(D)}_S,[f^{(D)}_S]
          \bigr).
        \]
  \item \textbf{law-level 部分 $L_S$：}
        取一个 law-level 系统 $L_S$，其 Jacobiator-gap
        为单一实数参数 $c_S>0$，并定义
        \[
          f^{(\JacGap)}_S(\gamma)
          =
          c_S\cdot f^{(D)}_S(\gamma),
          \qquad
          \gamma\in\Gamma_S.
        \]
        这给出偏好关系
        $\succ^{(\JacGap)}_S$，与 $\succ^{(D)}_S$ 显然相同。
        相应的路径语义对象为
        \[
          F_{\JacGap}(S)
          =
          \bigl(
            \Gamma_S,\succ^{(\JacGap)}_S,[f^{(\JacGap)}_S]
          \bigr).
        \]
\end{itemize}

定义 $\iota_S:F_D(S)\to F_{\JacGap}(S)$ 的底层映射为
$\mathrm{id}_{\Gamma_S}$，对应的重标度函数为
$\phi_{\iota_S}(x)=c_S x$。这显然是 $\mathsf{PathSem}$ 中的一个同构，
因而 $S$ 满足系统对象的定义。

\paragraph{同伦扭曲后的系统对象 $S'$}

现在考虑一次“同伦扭曲”之后的系统对象 $S'$，其路径集合保持不变：
$\Gamma_{S'}=\Gamma_S$。

\begin{itemize}
  \item \textbf{D 结构部分 $D_{S'}$：}
        反映“超限递归更新”后的新基准与逻辑性度量，取
        \[
          f^{(D)}_{S'}(\gamma)
          =
          \alpha\cdot f^{(D)}_S(\gamma),
          \qquad
          \alpha>0.
        \]
        这仍然诱导同一偏好关系
        $\succ^{(D)}_{S'}=\succ^{(D)}_S$，但在尺度上发生了
        严格单调重标度。
  \item \textbf{law-level 部分 $L_{S'}$：}
        设其 Jacobiator-gap 参数为 $c_{S'}>0$，并定义
        \[
          f^{(\JacGap)}_{S'}(\gamma)
          =
          c_{S'}\cdot f^{(D)}_{S'}(\gamma)
          =
          \alpha c_{S'}\cdot f^{(D)}_S(\gamma).
        \]
        同样诱导与 $\succ^{(D)}_S$ 一致的偏好关系
        $\succ^{(\JacGap)}_{S'}$。
\end{itemize}

只要取 $c_{S'}$ 满足
\[
  \alpha c_{S'} = c_S,
\]
就有
$f^{(\JacGap)}_{S'}(\gamma)=f^{(\JacGap)}_S(\gamma)$ 对所有路径成立。
此时，$F_{\JacGap}(S')$ 与 $F_{\JacGap}(S)$ 在 $\mathsf{PathSem}$
中实际上是同一个对象，而
$F_D(S')$ 则是 $F_D(S)$ 的一个严格单调重标度版本。
定义 $\iota_{S'}:F_D(S')\to F_{\JacGap}(S')$ 的底层映射为
$\mathrm{id}_{\Gamma_{S'}}$，对应的重标度函数为
$\phi_{\iota_{S'}}(x)=c_{S'} x$，即可得到系统对象 $S'$。

\paragraph{系统态射 $\Phi:S\to S'$}

构造映射
\[
  T_\Phi:\Gamma_S\longrightarrow\Gamma_{S'},
  \qquad
  T_\Phi=\mathrm{id}_{\Gamma_S},
\]
并令：
\begin{itemize}
  \item $F_D(\Phi)$ 的底层映射为 $T_\Phi$，
        对应重标度函数为
        $\phi_{F_D(\Phi)}(x)=\alpha x$，即
        $f^{(D)}_{S'}\circ T_\Phi=\phi_{F_D(\Phi)}\circ f^{(D)}_S$；
  \item $F_{\JacGap}(\Phi)$ 的底层映射为同一 $T_\Phi$，
        对应重标度函数取为恒等
        $\phi_{F_{\JacGap}(\Phi)}(x)=x$，
        因为在上面的参数选择下已经有
        $f^{(\JacGap)}_{S'}\circ T_\Phi=f^{(\JacGap)}_S$。
\end{itemize}

这样，(SM1) 与 (SM2) 显然满足；而 (SM3) 要求
\[
  F_{\JacGap}(\Phi)\circ\iota_S
  =
  \iota_{S'}\circ F_D(\Phi).
\]
在评分函数层面，这对应于严格单调重标度函数之间的关系
\[
  \phi_{F_{\JacGap}(\Phi)}\circ\phi_{\iota_S}
  =
  \phi_{\iota_{S'}}\circ\phi_{F_D(\Phi)},
\]
代入
$\phi_{F_{\JacGap}(\Phi)}=\mathrm{id}$、
$\phi_{\iota_S}(x)=c_S x$、
$\phi_{\iota_{S'}}(x)=c_{S'} x$、
$\phi_{F_D(\Phi)}(x)=\alpha x$，
可得
\[
  c_S x = c_{S'}\alpha x,
\]
这正是前面对 $c_{S'}$ 的约束条件。因此，上述构造给出了
一个满足所有公理的系统态射 $\Phi:S\to S'$。

从语义上看，$S$ 表示某一时刻下的 D 结构 / law-level 配置，
而 $S'$ 表示沿着一次“同伦扭曲 + 超限递归更新”之后的新配置。
系统态射 $\Phi$ 把同一条路径在这两个配置之间进行匹配，
并通过严格单调重标度保持 D 语义与 Jacobiator-gap 语义的一致性，
因此是对前文“同伦扭曲语义等价性”在 $\mathsf{Sys}$ 中的一个
具体模型。

\section{从 D 结构到 Jacobiator-gap：构造方向一}

\subsection{构造 Jacobiator-gap 兼容的 law-level 系统}

给定一个固定的 D 结构状态 $(E,H,w)$。
在该状态下，我们考虑所有可行的候选路径集合 $\Gamma(E,H)$，
以及对应的逻辑性度量
$\mathcal{L}_D(\gamma;w)$。

在 G-Framework 的语言中，我们可以做如下构造：
\begin{enumerate}
  \item 在 law-space $\mathcal{L}$ 中，构造一个 law-level 系统
        $L_w$，其低阶观测量（例如期望收益、安全度量等）与
        $\mathcal{L}_D$ 的各个分量 $d_i(\gamma)$ 一一对应；
  \item 将 D 结构的“基准权重向量” $w$ 编码为 $L_w$ 的参数，
        使得 $L_w$ 对路径 $\gamma$ 的 law-space 评价可以写成
        \[
          F\bigl(\JacGap(L_w),\gamma\bigr)
          = \psi_w\Bigl(\mathcal{L}_D(\gamma;w)\Bigr),
        \]
        其中 $\psi_w$ 是某个严格单调递增函数。
\end{enumerate}

这一构造在主稿的框架下，可以理解为：
在 law-space 中寻找一个同伦 PFB-GNLA 结构 $L_w$，使得
Jacobiator-gap 证书 $\JacGap(L_w)$ 在数值上精确编码了
D 结构的逻辑性度量。

\subsection{语义层面的单调同构}

在上述构造下，对任意两条路径 $\gamma_1,\gamma_2$，有：
\[
  \mathcal{L}_D(\gamma_1;w)>\mathcal{L}_D(\gamma_2;w)
  \;\Longleftrightarrow\;
  \psi_w\Bigl(\mathcal{L}_D(\gamma_1;w)\Bigr)
  >
  \psi_w\Bigl(\mathcal{L}_D(\gamma_2;w)\Bigr),
\]
又由于
\[
  \psi_w\Bigl(\mathcal{L}_D(\gamma;w)\Bigr)
  =
  F\bigl(\JacGap(L_w),\gamma\bigr),
\]
可得
\[
  \gamma_1\succ_D\gamma_2
  \;\Longleftrightarrow\;
  \gamma_1\succ_{\JacGap}\gamma_2.
\]

因此，在评分函数之间取
$
  f_1(\gamma)=\mathcal{L}_D(\gamma;w),
$
$
  f_2(\gamma)=F\bigl(\JacGap(L_w),\gamma\bigr),
$
并令 $\phi=\psi_w$，就得到 $\succ_D$ 与 $\succ_{\JacGap}$ 在
路径偏好意义下语义等价。

\section{从 Jacobiator-gap 到 D 结构：构造方向二}

\subsection{由证书型不变量回构逻辑性度量}

反过来，给定一个 law-level 系统 $L$ 及其 Jacobiator-gap 证书
$\JacGap(L)$，我们可以在 O3 理论的语境下构造一个 D 结构，步骤为：
\begin{enumerate}
  \item 把 $\JacGap(L)$ 看作是一组“结构性约束”的向量，
        例如对安全性、鲁棒性、一致性等指标的下界约束；
  \item 在 D 结构中引入维度分量 $d_i(\gamma)$，分别对应这些约束在路径
        $\gamma$ 上的表现，例如“违反 Jacobi 恒等的程度”导致的
        风险增量；
  \item 令 D 结构的权重向量 $w$ 由 Jacobiator-gap 的数值决定，
        例如
        \[
          w_i = g_i\bigl(c_i(L)\bigr),
        \]
        其中 $c_i(L)$ 是 Jacobiator-gap 的分量，$g_i$ 是严格单调的
        正函数；
  \item 定义逻辑性度量算子
        \[
          \mathcal{L}_D(\gamma;w)
          =
          \sum_i w_i\cdot d_i(\gamma),
        \]
        并选取 $d_i(\gamma)$ 使得
        \[
          \mathcal{L}_D(\gamma;w)
          =
          \tilde{\psi}\Bigl(
            F\bigl(\JacGap(L),\gamma\bigr)
          \Bigr),
        \]
        其中 $\tilde{\psi}$ 是严格单调递增函数。
\end{enumerate}

这一步可以理解为：把 Jacobiator-gap 视作“结构性风险证书”，再把它
翻译成 D 结构内部的“逻辑性度量权重”，使得 D 结构在路径上的评价
与 Jacobiator-gap 导出的评价同构。

\subsection{反向语义等价性的证明}

在上述构造下，对任意路径 $\gamma_1,\gamma_2$ 有：
\[
  F\bigl(\JacGap(L),\gamma_1\bigr)
  >
  F\bigl(\JacGap(L),\gamma_2\bigr)
  \;\Longleftrightarrow\;
  \tilde{\psi}\Bigl(
    F\bigl(\JacGap(L),\gamma_1\bigr)
  \Bigr)
  >
  \tilde{\psi}\Bigl(
    F\bigl(\JacGap(L),\gamma_2\bigr)
  \Bigr),
\]
又因为
\[
  \tilde{\psi}\Bigl(
    F\bigl(\JacGap(L),\gamma\bigr)
  \Bigr)
  =
  \mathcal{L}_D(\gamma;w),
\]
所以
\[
  \gamma_1\succ_{\JacGap}\gamma_2
  \;\Longleftrightarrow\;
  \gamma_1\succ_D\gamma_2.
\]

从而，在评分函数之间取
$
  f_1(\gamma)=F\bigl(\JacGap(L),\gamma\bigr),
$
$
  f_2(\gamma)=\mathcal{L}_D(\gamma;w),
$
并令 $\phi=\tilde{\psi}$，即可得到 $\succ_{\JacGap}$ 与 $\succ_D$
在路径偏好意义下语义等价。

\section{泛拓扑与性变态射视角下的补充证明}

\subsection{kernel\_reference 中的泛拓扑与泛抽象代数框架}

在文献 \cite{KernelReferenceGenericTopology,KernelReferenceSexualOperators} 中，已经给出了一个面向数学结构的\emph{泛拓扑（generic topology）}
与\emph{泛抽象代数}的统一框架，并在其中引入了
性变算子 $T$ 与性变态射 $f$：
\[
  T:(\mathcal{S},P)\to(\mathcal{S}',P'),
  \qquad
  f:(\mathcal{S},P)\to(\mathcal{S}',P'),
\]
其中 $\mathcal{S}$、$\mathcal{S}'$ 是数学结构，
$P$、$P'$ 是性质集合。

为了在后文中更精确地使用“共性（交集）拓扑”的语言，我们在此给出
一个最小形式化模型。记 $\mathcal{P}_{\mathrm{all}}$ 为允许的
“原子性质”的全集，每个性质集合 $P$ 视为 $\mathcal{P}_{\mathrm{all}}$
的子集。定义性质空间
\[
  \mathsf{Top}_P
  =
  \bigl\{\,P\subseteq\mathcal{P}_{\mathrm{all}}\ \bigm|\ 
    P \text{ 满足 kernel\_reference 中给定的基本一致性约束}\,\bigr\},
\]
并在 $\mathsf{Top}_P$ 上引入如下拓扑基：
\[
  \mathcal{B}
  =
  \bigl\{
    U(F)
    \ \bigm|\ 
    F\subseteq\mathcal{P}_{\mathrm{all}}
    \text{ 有限},\ 
    U(F)=\{\,P\in\mathsf{Top}_P\mid F\subseteq P\,\}
  \bigr\}.
\]
易见 $\mathcal{B}$ 在有限交下封闭，并生成一个拓扑 $\tau_P$；
我们把拓扑空间 $(\mathsf{Top}_P,\tau_P)$ 称为“共性（交集）拓扑”，
因为每个基开集 $U(F)$ 精确对应一组同时被满足的共性性质。

在时间参数空间 $I=[0,1]$ 上取标准实数拓扑。给定一族性质集合
$\{P_t\}_{t\in I}$，若映射
\[
  \gamma_P:I\longrightarrow\mathsf{Top}_P,\quad
  \gamma_P(t)=P_t
\]
在上述拓扑下连续，则称 $\{P_t\}$ 为一条\emph{性质路径}。
“保持共性性质不突变”可以形式化为：
存在有限性质子集 $F\subseteq\mathcal{P}_{\mathrm{all}}$，
使得对所有 $t\in I$ 都有 $F\subseteq P_t$，即
$\gamma_P(I)\subseteq U(F)$。

按照上述文档的设定，可以把这一框架理解为：
\begin{itemize}
  \item \textbf{泛拓扑层：}
        把“具有某种性质的结构”视为一个点，把“满足同一性质子族的结构族”
        视为一个“开集”，从而在性质空间上诱导出一种\emph{含交集的拓扑}：
        交集对应“共性（交集）性质”的同时满足；
  \item \textbf{泛抽象代数层：}
        在每个性质子族内部，给出相应的泛抽象代数结构
        （运算、关系与约束），并用性变算子 $T$ 描述
        性质集 $P\to P'$ 演化时代数规则的变化；
  \item \textbf{路径层：}
        用性变态射 $f$ 描述在泛拓扑下沿某条“性质路径”的演化，
        即 $(\mathcal{S},P)\to(\mathcal{S}',P')$ 的拓扑路径变换。
\end{itemize}

从这个视角看，性变算子与性变态射本身就已经把“拓扑—代数—路径”
三层结构统一在了同一个元数学框架下。

\subsection{性变态射（性变算子）与传统拓扑同伦的对价}

我们在此给出一个补充性命题，说明在上述泛拓扑框架中，
性变态射（及其伴随的性变算子）如何\emph{价于}传统数学中的拓扑同伦。

\begin{proposition}[交集拓扑下的同伦类对价]
在 kernel\_reference 的设定与上述 $(\mathsf{Top}_P,\tau_P)$ 构造下，
考虑所有以性质集为指标构造的泛拓扑空间 $\mathsf{Top}_P$，
其开集由共性性质的交集生成。
则：
  \begin{enumerate}
    \item 每一条由性变态射族
          $\{f_t:(\mathcal{S},P_t)\to(\mathcal{S}_t,P_t)\}_{t\in[0,1]}$
          所定义的性质路径 $\{P_t\}$，若 $\gamma_P:I\to\mathsf{Top}_P$
          连续，则对应于 $\mathsf{Top}_P$ 中一条传统意义上的路径；
    \item 反之，每一条在 $\mathsf{Top}_P$ 中保持共性性质不突变的
          连续路径 $\gamma_P$，都可以被提升为一族性变态射
          $\{f_t\}$，使得 $f_t$ 在每个时刻 $t$ 上尊重
          kernel\_reference 中给出的性质偏序约束；
    \item 在自然的等价关系下（两族性变态射若诱导同伦等价的
          $\gamma_P$ 则视为等价），性变态射族的等价类与
          含交集拓扑 $\mathsf{Top}_P$ 中的同伦类一一对应。
  \end{enumerate}
\end{proposition}

\begin{proof}[证明思路]
我们只给出与本笔记目的相关的证明要点：
\begin{itemize}
  \item 由 kernel\_reference 中对性变态射的定义可知，
        $f:(\mathcal{S},P)\to(\mathcal{S}',P')$ 必须遵守
        性质集合上的偏序与路径一致性；
        因此沿着参数 $t$ 的性变态射族
        $\{f_t\}_{t\in[0,1]}$，自然诱导出性质路径
        $\{P_t\}$ 以及映射 $\gamma_P:I\to\mathsf{Top}_P$；
        当 $\gamma_P$ 连续时，它就是拓扑意义下的一条路径；
  \item 含交集拓扑 $\mathsf{Top}_P$ 的开集由性质子族的交集生成，
        这保证了只要性质路径在 $t$ 上不发生“突变跳跃”，
        即对任意有限 $F$ 满足 $F\subseteq P_{t_0}$ 时，
        在 $t_0$ 附近仍有 $F\subseteq P_t$，就可以在传统意义下
        视为一条连续路径，并进一步讨论同伦；
  \item 另一方面，给定 $\mathsf{Top}_P$ 中一条同伦路径，
        我们可以在每个时间片 $t$ 上选取与之相容的结构
        $(\mathcal{S}_t,P_t)$，并利用 kernel\_reference 中的
        泛抽象代数构造相应的性变态射 $f_t$，
        从而把拓扑同伦提升为“性质+结构”的动态演化；
  \item 在自然的等价类（例如两族性变态射若诱导 $\gamma_P$ 之间的
        同伦等价则视为等价）上，上述双向构造互为逆，从而建立
        “一一对应”的结论。
\end{itemize}
因此，从语义上讲，性变态射族所定义的演化，与传统数学中在
共性（交集）拓扑下的同伦类是等价的，只是前者在 kernel\_reference
框架中显式追踪了性质集与代数结构的同步演化。
\end{proof}

\subsection{上同调类、雅可比间隙与泛拓扑障碍}

在 G-Framework 的纯数学主稿中，雅可比间隙 $\JacGap$ 可以被理解为
一种“law-level 的上同调障碍”：当某个 law-level 系统 $L$ 的
Jacobi 恒等式无法在全局上被严格实现时，其局部修正的失配部分
可以被编码到某个上同调类中，记为
\[
  [\JacGap(L)]\in H^{k+1}(X;\Lambda),
\]
其中 $X$ 是适当的基空间，$\Lambda$ 是目标系数系。

在 kernel\_reference 的泛拓扑框架中，可以做出如下对应性解释：
\begin{itemize}
  \item 性质空间 $P$ 及其共性拓扑 $\mathsf{Top}_P$ 中，
        若存在某些性质组合在全局上无法被一致实现，
        则这些“无法一致实现的性质环”构成了一个上同调障碍；
  \item 该障碍意味着：无论怎样选取性变态射族 $\{f_t\}$，
        都无法在不引入“性质断裂”的前提下，把某一类结构连续地
        变形到另一类结构；
  \item 这与 G-Framework 中 Jacobiator-gap 所刻画的情形完全同构：
        当 $\JacGap(L)\neq 0$ 时，law-space 中存在某种“扭曲”，
        阻止我们在保持低阶观测量不变的前提下把 $L$ 提升为
        严格 Jacobi 的对象。
\end{itemize}

因此，可以把
\[
  \text{“上同调类”} \quad\Longleftrightarrow\quad \text{“雅可比间隙”}
\]
理解为：泛拓扑下性质路径的\emph{障碍或断路}的抽象指标。
一旦该上同调类非平凡，就意味着存在某种“不可消去的性质环”，
使得性质演化不可在有限同伦内实现全局连续化。

\subsection{D 结构与同伦扭曲的超限递归作用}

kernel\_reference 中关于 D 结构、DERI 算法及性变算子/性变态射的讨论，
为我们提供了一个关键观点：\emph{D 结构可以被视为在泛拓扑/泛代数之上
执行“超限递归”的逻辑引擎}。

结合前文，我们可以做如下解释：
\begin{itemize}
  \item 当上同调类 $[\JacGap(L)]$ 非平凡时，在单一拓扑尺度上，
        性变态射族无法消去该“性质环”，从而在 $\mathsf{Top}_P$ 中
        形成“拓扑断路”；
  \item D 结构通过逻辑性度量 $\mathcal{L}_D$ 与环境/历史 $(E,H)$ 的
        自反性更新，可以在“更高一层”的元空间上迭代调整基准 $w$，
        相当于对 law-space 进行“同伦扭曲”（homotopy twist）：
        不直接在原来的同伦类内强行填补空洞，而是通过改变
        “允许的性质路径族”来绕过该上同调类；
  \item 若把这一过程视为对“路径空间”本身的迭代，则可以用
        超限递归的语言来刻画：每一次 D 结构更新都是在更高阶的
        性质空间上添加一个新的“层级”，从而在极限过程中
        实现对原有上同调障碍的“间接消解”。
\end{itemize}

从泛拓扑视角看，这等价于：在原有 $\mathsf{Top}_P$ 上添加
一个“扩展拓扑” $\widetilde{\mathsf{Top}}_P$，其开集不仅记录
性质交集，还记录“允许的超限迭代路径”。在这一扩展后的拓扑中，
原先的断路被绕开，从而实现“拓扑继续”。

\subsubsection*{一个极简形式模型示例}

为了在语言上捕捉上述“超限递归 + 同伦扭曲”机制，我们给出一个
极简形式模型，仅作为今后形式化工作的参考：
\begin{itemize}
  \item 取按序数索引的一族 D 结构
        $\{D_\alpha=(E_\alpha,H_\alpha,w_\alpha,\mathcal{L}_\alpha)\}
        _{\alpha<\kappa}$，其中 $\kappa$ 是某个极限序数；
  \item 对每个 $\alpha$，令 $\Gamma_\alpha$ 为在 D 结构 $D_\alpha$
        下可行的路径集合，并令
        \[
          f_\alpha:\Gamma_\alpha\longrightarrow\mathbb{R}
        \]
        为与 $\succ_{D_\alpha}$ 相容的评分函数；
  \item 定义递推：
        \[
          D_{\alpha+1}
          =
          \mathsf{Update}\bigl(
            D_\alpha,
            \JacGap(L_\alpha)
          \bigr),
        \]
        其中 $L_\alpha$ 是与 $D_\alpha$ 对应的 law-level 系统，
        $\mathsf{Update}$ 把当前无法消去的雅可比间隙编码进
        新的基准 $w_{\alpha+1}$ 与新的可行路径集 $\Gamma_{\alpha+1}$；
  \item 在极限序数 $\lambda<\kappa$ 处，令
        \[
          \Gamma_\lambda
          =
          \bigcup_{\alpha<\lambda}\Gamma_\alpha,
          \qquad
          f_\lambda(\gamma)
          =
          \sup_{\alpha<\lambda} f_\alpha(\gamma|_{\Gamma_\alpha}),
        \]
        并由 $f_\lambda$ 诱导新的偏好关系 $\succ_{D_\lambda}$。
\end{itemize}

在这个模型中，每一步 $\alpha\mapsto\alpha+1$ 都可以被视为对
路径空间同伦类型的一次“扭曲”，目的是在保持低阶可接受性的前提下
扩展可行路径集，从而逐步绕开由 $\JacGap$ 所给出的上同调障碍。
在极限阶段 $\lambda$，我们得到一个汇总了所有先前扭曲结果的
“极限 D 结构” $D_\lambda$，对应着在扩展拓扑
$\widetilde{\mathsf{Top}}_P$ 中实现的拓扑继续。

\paragraph{在无限同伦与主丛联络下刻画修复深度与广度}

为与 G-Framework 主稿中 law-space 的主丛几何与三层联络严格对标，
我们在上述极简模型的基础上引入一个主丛版本的统一抽象，并显式
采用主稿中的\emph{同伦 PFB-GNLA 版本}。设
\[
  \pi:P\longrightarrow X
\]
是一个具有结构群 $G$ 的主丛。

\begin{definition}[同伦 PFB-GNLA law-level 结构]
称
\[
  L=(P,\omega,H)
\]
是一个同伦 PFB-GNLA law-level 结构，如果：
  \begin{itemize}
    \item $\omega\in\Omega^1\bigl(P;\widehat{\mathfrak g}\bigr)$ 是
          一个 $G$-等变的主丛联络 1-形式，取值于某个扩展的
          $\widehat{\mathfrak g}$（其零级部分为 $\mathfrak g$），
          曲率为
          \[
            F_\omega
            =
            \mathrm{d}\omega
            + \tfrac{1}{2}[\omega,\omega]；
          \]
    \item 在 $\operatorname{ad}(P)$ 或其扩展上给定一个 2-项
          $L_\infty$-代数 $(l_1,l_2,l_3)$，其中 $l_2$ 是
          Lie 括号，$l_3$ 的系数由一个闭 3-形式
          $H\in\Omega^3(X;\Lambda)$ 控制；
    \item Jacobi 恒等式的失败由
          \[
            l_2\bigl(l_2(a,b),c\bigr)+\text{循环置换}
            =
            l_1\bigl(l_3(a,b,c)\bigr)
          \]
          及 $H$ 满足的 Bianchi 型方程统一编码，并且 Jacobiator-gap
          的上同调类可以取为
          \[
            [\JacGap(L)]
            =
            [H]\in H^3(X;\Lambda),
          \]
          更高阶情形（$H^{k+1}(X;\Lambda)$）可类似处理。
  \end{itemize}
\end{definition}

在这一设定下，我们可以把“雅可比间隙”理解为 $H$ 的上同调类所刻画的
同伦 PFB-GNLA 障碍。为简洁起见，下文仍记
\[
  [\JacGap(L)]\in H^{k+1}\bigl(X;\Lambda\bigr)
\]
来表示相应的上同调指标，其中系数系 $\Lambda$ 在典型情形下可以具体化为与
$\operatorname{ad}P$ 相关的某个表示。按照本笔记前文的双向构造，
前述按序数索引的 D 结构族
$\{D_\alpha\}_{\alpha<\kappa}$ 唯一对应于一族 law-level 结构
\[
  L_\alpha=(P,\omega_\alpha,H_\alpha),\qquad \alpha<\kappa,
\]
其中每个 $\omega_\alpha$ 是 $P$ 上的一个 PFB-GNLA 联络，
满足 $\JacGap(L_\alpha)$ 与 $D_\alpha$ 的路径语义等价。

\begin{definition}[同伦扭曲深度与广度的主丛版本]
称 $\{L_\alpha=(P,\omega_\alpha)\}_{\alpha<\kappa}$ 是 $L_0$ 的一条
\emph{无限同伦扭曲塔}，如果对每个 $\alpha<\kappa$，存在一个
主丛意义下的同伦
\[
  H_\alpha:P\times I\longrightarrow P,
\]
满足：
  \begin{itemize}
    \item $H_\alpha$ 为 $G$-等变并覆盖恒等映射
          $\mathrm{id}_X$；
    \item $H_\alpha(\cdot,0)$ 与 $H_\alpha(\cdot,1)$ 分别对应联络
          $\omega_\alpha$ 与 $\omega_{\alpha+1}$（即同伦两端的
          水平分布分别给出这两个 law-level 结构）；
    \item 同伦路径 $\{H_\alpha(\cdot,t)\}_{t\in I}$ 始终停留在主稿中
          允许的 PFB-GNLA 结构类中。
  \end{itemize}
在这样的无限同伦扭曲塔上，我们定义：
  \begin{itemize}
    \item \textbf{同伦扭曲深度}为
          \[
            \alpha_\ast
            =
            \min\bigl\{\alpha<\kappa \mid
              [\JacGap(L_\alpha)]=0\ \text{（或落入预设容忍子集）}
            \bigr\},
          \]
          即沿着主丛同伦塔首次在上同调意义下消解（或弱化到容忍阈值）
          雅可比间隙所需要的最小序数高度；
    \item 记
          \[
            J_\ast
            =
            \bigl\{\,i \,\bigm|\,
              \exists \alpha<\alpha_\ast,\ 
              \text{在某个同伦 }H_\alpha\text{ 中，}
              c_i\bigl(L_\alpha\bigr)
              \text{发生非平凡的同伦更新}
            \bigr\},
          \]
          其中 $\JacGap(L_\alpha)=(c_1(L_\alpha),\dots,c_k(L_\alpha))$，
          则称 $|J_\ast|$ 以及 $J_\ast$ 在系数系 $\Lambda$ 中张成的
          最小子对象为该修复过程的\emph{同伦扭曲广度}：
          它刻画了需要在多少个独立的 law-level 方向上展开扭曲。
  \end{itemize}
\end{definition}

\begin{remark}[连续统崩溃$\Rightarrow$深度，同伦扭曲$\Rightarrow$广度（GZ-CBT 语义下的耦合表述）]
在本笔记的语言里，“雅可比间隙的深度/广度”默认指\emph{修复深度/修复广度}：
深度用 $\alpha_\ast$（或主稿中的 $n_\ast,n_{\mathrm{tower}},d_{\mathrm{sector}}$）刻画
“要沿着修复链/塔走多深才能把 JacGap 消去或压到容忍阈值”；广度用 $J_\ast$
刻画“修复过程中需要同时动用多少个独立的证书分量/方向”。

把这一语言对齐到 \cite{GaoZheng2025GFramework} 的 GZ-CBT 叙事，
可以更精确地陈述为：
\begin{itemize}
  \item \textbf{连续统崩溃（CH breakdown）通常推高修复深度：}
        主稿把连续统侧的“崩溃/失配”量化为一个时间局部指标（符号略同）：
        \[
          \mathcal{B}_{\mathrm{CH}}(t)
          =
          w_{\Delta}\,\Delta_{\mathrm{CH}}(t)
          +
          w_{B}\,B_{\mathrm{CH}}(t),
        \]
        并与算子侧的 Jacobiator-gap 指标
        \[
          J_{\mathrm{op}}(t)
          :=
          \bigl\|
            \JacGap\bigl(L^{\mathrm{op}}_{\mathrm{cont}}(t)\bigr)
          \bigr\|
        \]
        合成为
        \[
          \mathcal{B}_{\mathrm{CH+op}}(t)
          =
          \mathcal{B}_{\mathrm{CH}}(t)
          +
          w_{\mathrm{op}}\,J_{\mathrm{op}}(t).
        \]
        CBT--OHU scheme 的“修复深度”在工程上体现为：把
        $\mathcal{B}_{\mathrm{CH+op}}$（或其能量泛函）压回阈值区间所需的最小步数
        $n_\ast$、最小塔阶数 $n_{\mathrm{tower}}$、或最小扇区迁移距离
        $d_{\mathrm{sector}}$。它们在本笔记中统一投影为序数刻度 $\alpha_\ast$。
        因而连续统崩溃越严重（$\mathcal{B}_{\mathrm{CH}}$ 越大），通常意味着要走得更“深”。
  \item \textbf{同伦扭曲通常推高修复广度：}
        在主稿的 homotopy chart / $H$-twist 叙事中，Jacobi 失败被显式化为
        $\ell_3^H$、更高括号与证书分量
        $\JacGap(L)=(c_1,\dots,c_k)$，并需要跨 layer/chart 协同调节。
        在本节的主丛语言里，这等价于\emph{扩展纤维方向/结构群方向}
        （例如 $\mathfrak g\leadsto\widehat{\mathfrak g}$，或把可见的证书方向做直和扩展），
        从而使得“参与非平凡更新的分量集合” $J_\ast$ 变大——这就是广度。
\end{itemize}

用两条要点把上述对位压缩为你关心的表达就是：
\begin{itemize}
  \item \textbf{连续统崩溃往往迫使我们显式引入（或启用）主丛/三层联络的几何包裹，从而体现为“向更高层级推进”的深度。}
  \item \textbf{同伦扭曲往往迫使我们在主丛包中扩展纤维/结构群的独立方向，从而体现为“需要同时处理多少方向”的广度。}
\end{itemize}

最后强调：以上是\emph{典型耦合关系}，不是严格的一一因果分解——
CH breakdown 也可能迫使进入 homotopy chart 从而“变宽”，同伦扭曲也可能需要多轮 OHU 才能“变深”。
因此把现象拆成“深度/广度”更多是一种\emph{坐标选择/认知组织的便利}：在不改变核心诊断量与修复语义的前提下，
你可以拆分，也可以不拆分；关键在于全篇口径一致、并明确哪些量是标量强度（如 $\mathcal{B}_{\mathrm{CH}}(t)$），哪些量是支持/方向信息（如 $S_{\mathrm{CH}}(L),B_{\mathrm{chart}}(L),J_\ast$）。
\end{remark}

\begin{proposition}[超限 D 结构序列与主丛无限同伦的对标]
在本笔记的公理化假设下，每一条满足前述递推式的 D 结构超限序列
  $\{D_\alpha\}_{\alpha<\kappa}$ 都唯一对应于一条主丛上的
  无限同伦扭曲塔 $\{L_\alpha=(P,\omega_\alpha)\}_{\alpha<\kappa}$，
  并且：
  \begin{enumerate}
    \item D 结构意义下以 $\{D_\alpha\}$ 定义的“雅可比间隙修复深度”
          恰好等于上式中的 $\alpha_\ast$；
    \item D 结构意义下定义的“雅可比间隙修复广度”
          （参与实际更新的 Jacobiator-gap 分量族）在自然同构意义下
          与 $J_\ast$ 一致。
  \end{enumerate}
\end{proposition}

\begin{proof}[证明思路]
由前文的语义等价定理，每个 D 结构 $D_\alpha$ 唯一对应于一个
Jacobiator-gap 证书体系 $(L_\alpha,\JacGap)$，从而对应于主丛上的
某个 PFB-GNLA 联络 $\omega_\alpha$。利用 law-space 中
PFB-GNLA 结构在同伦意义上的完备性，可以在每一对相邻的
$\omega_\alpha$ 与 $\omega_{\alpha+1}$ 之间选取一条保持结构约束的
主丛同伦 $H_\alpha$，从而在主丛几何上重建出一条
无限同伦扭曲塔 $\{L_\alpha\}_{\alpha<\kappa}$。

另一方面，$\mathsf{Update}(D_\alpha,\JacGap(L_\alpha))$ 的构造方式
保证：在从 $\alpha$ 到 $\alpha+1$ 的更新中，恰好是那些
$c_i(L_\alpha)$ 被用于修改基准 $w_{\alpha+1}$ 与可行路径集
$\Gamma_{\alpha+1}$。因此，D 结构视角下“真正参与修复”的分量指标集
与上面定义的 $J_\ast$ 一致；而“修复完成”的判据就是
$[\JacGap(L_\alpha)]$ 在上同调意义下进入零类或预设容忍子集中，
于是给出的临界序数与主丛版本的 $\alpha_\ast$ 相同。
综上，D 结构视角下定义的“同伦扭曲深度/广度”与基于主丛联络的
无限同伦刻度严格对标。
\end{proof}

\begin{definition}[GZ-TLC 型 law-space 主丛与 law-level 投影]
设 $\widehat{P}\to X_{\mathrm{tot}}$ 是一个主丛，其中
\[
  X_{\mathrm{tot}}
  =
  X\times\Lambda_{\mathrm{par}}\times\mathsf{GZLS},
\]
$X$ 对应系统/环境层，$\Lambda_{\mathrm{par}}$ 为参数层，
$\mathsf{GZLS}$ 为 law-space 层。称
\[
  \mathcal{A}_{\mathrm{GZTLC}}
  \in
  \Omega^1\bigl(\widehat{P};\widehat{\mathfrak g}\bigr)
\]
为一条 GZ-TLC 型三层总联络，如果其曲率
  \[
    \mathcal{F}_{\mathrm{GZTLC}}
    =
    \mathrm{d}\mathcal{A}_{\mathrm{GZTLC}}
    + \tfrac{1}{2}
      [\mathcal{A}_{\mathrm{GZTLC}},\mathcal{A}_{\mathrm{GZTLC}}]
  \]
可以按“几何—参数—法则”三层分解，且存在一个嵌入
\[
  \iota:P\hookrightarrow\widehat{P},
\]
使得某个\emph{law-level 分量} $\mathcal{A}^{\mathrm{law}}$ 满足
  \[
    \omega
    =
    \iota^\ast\mathcal{A}^{\mathrm{law}}
  \]
是主丛 $P$ 上的同伦 PFB-GNLA 联络，从而给出一个 law-level 结构
$L=(P,\omega,H)$。
\end{definition}

在这一意义下，每个 law-level 结构 $L_\alpha$ 都可以视为某条
GZ-TLC 型总联络在 law-level 的投影；沿着 GZ-TLC 水平的同伦扭曲，
可以在不离开 PFB-GNLA 类的前提下演化 $L_\alpha$。

\begin{theorem}[雅可比间隙修复的主丛同伦--GZ-TLC 对标]
在假设 \textnormal{(A1)}--\textnormal{(A3)} 以及上述同伦 PFB-GNLA 与
  GZ-TLC 型 law-space 主丛设定下，固定一个初始 law-level 结构
  $L_0=(P,\omega_0,H_0)$。则以下两类数据在自然的同伦--规范等价意义下
  一一对应：
  \begin{enumerate}
    \item 一条满足递推
          $D_{\alpha+1}
           =\mathsf{Update}\bigl(D_\alpha,\JacGap(L_\alpha)\bigr)$
          的 D 结构超限序列 $\{D_\alpha\}_{\alpha<\kappa}$，其
          雅可比间隙修复深度与广度记为 $(\alpha_\ast,J_\ast)$；
    \item 一条 GZ-TLC 总联络
          $\{\mathcal{A}_{\mathrm{GZTLC},\alpha}\}_{\alpha<\kappa}$
          的无限同伦扭曲塔，其 law-level 分量
          \[
            \omega_\alpha
            =
            \iota^\ast
            \mathcal{A}^{\mathrm{law}}_{\mathrm{GZTLC},\alpha}
          \]
          给出 law-level 结构 $L_\alpha=(P,\omega_\alpha,H_\alpha)$，
          并在前述主丛意义下具有同样的修复深度与广度
          $(\alpha_\ast,J_\ast)$。
  \end{enumerate}
换言之，“雅可比间隙的修复，包括同伦扭曲深度和广度”既可以在
  D 结构--路径语义侧刻画，也可以在 GZ-TLC 型主丛--无限同伦侧
  严格等价地重述。
\end{theorem}

\begin{proof}[证明思路]
命题关于“超限 D 结构序列与主丛无限同伦的对标”已经给出了
$\{D_\alpha\}$ 与 $\{L_\alpha\}$ 之间的一一对应，并证明了两个侧面
的修复深度/广度刻度一致。另一方面，同伦 PFB-GNLA 的定义保证
每个 $L_\alpha$ 可以被视为某条 GZ-TLC 总联络的 law-level 投影，
而 GZ-TLC 的三层结构提供了从几何/参数扭曲到 law-level 扭曲的
标准提升方式。沿着一条 GZ-TLC 水平的无限同伦扭曲塔
$\{\mathcal{A}_{\mathrm{GZTLC},\alpha}\}$，其 law-level 分量族
$\{\omega_\alpha\}$ 正好重现了先前构造的
$\{L_\alpha\}=(P,\omega_\alpha,H_\alpha)$，从而两类数据在自然的
同伦--规范等价下互为逆构造。

反之，从任一条 GZ-TLC 型无限同伦扭曲塔出发，通过 law-level 投影
即可恢复 $\{L_\alpha\}$，再利用路径语义等价与
Update 递推构造出对应的 $\{D_\alpha\}$。由于 Jacobiator-gap
通过同伦 PFB-GNLA 的闭 3-形式（或更高阶闭 $(k+1)$-形式）编码，
其修复深度与广度在两个视角下给出的正是同一个上同调与分量簇数据，
从而定理成立。
\end{proof}

\begin{remark}[与主稿中 Jacobiator-gap 刻度的关系]
结合 \cite{GaoZheng2025GFramework} 可以更精确地理解
“扭曲深度/广度”与主稿中现有结构的对齐方式：
\begin{itemize}
  \item \textbf{广度刻度：}
        主稿中的 CH 分层（$l_2$、$l_{\ge3}$、一般的 $l_n$），
        再加上四个 PFB-GNLA chart
        （strict / homotopy /
        variable functional--operator / original law-level），
        本身就刻画了“有多少个语义/同伦层被激活”。
        这与本笔记所说的“扭曲广度（涉及多少层/多少方向）”是对得上的，
        只是主稿采用的是 layer 与 chart 的语言，
        没有单独命名为“广度”。
  \item \textbf{深度 $\approx$ 重写链长 / law-bundle 阶数 / 扇区迁移：}
        主稿中还出现了 CBT--OHU 重写序列的步数 $n$ 及其
        JacGap-minimal 极限点、law-bundle 塔
        $\pi_n:\mathcal{E}_n\to B$ 的阶数 $n$，以及在 GZ-LS 中按
        CH-energy/JacGap 把 law-space 划分为不同“区域/扇区”并讨论
        其间的重写与流动。这些结构在事实上传递的是
        “我们在多长的一条链、多深的一层上处理扭曲”的信息，
        可以解释为一种“扭曲深度”（扇区内的进深、跨扇区的层数），
        只是主稿目前使用的是“步数、阶数、区域/sector”的表述，
        并未显式抽象为形式化的不变量。
\end{itemize}

综合来看：
\begin{itemize}
  \item 主稿已经具备横向的“广度”刻度：通过 $l_n$、chart、
        law-bundle 层级说明在哪些层/方向上观察 JacGap；
  \item 也具备纵向的“深度”刻度：通过重写步数、塔的阶数、
        不同 CH/CBT 扇区之间的迁移来度量“扭曲走了多远”。
\end{itemize}

本笔记在此基础上，引入“同伦扭曲深度（序数高度）/广度
（参与分量子空间）”以及与 GZ-TLC 型主丛的严格对标，一方面是对
主稿已有结构的统一重述，另一方面把原本分散在不同章节和语境中的
刻度收束成一个明确的“扭曲深度/广度”语言，并与 D 结构、
主丛无限同伦建立了显式的一一对应关系。
\end{remark}

\section{G-Framework 主稿中的扭曲广度与扭曲深度刻度}

在前一节中，我们从 D 结构与同伦 PFB-GNLA 的视角引入了“扭曲广度”
与“扭曲深度”的刻度。本节转而站在
\cite{GaoZheng2025GFramework} 的内部坐标系中，精确解释主稿中
已有的 CH-layering、PFB-GNLA 四图表、law-bundle 塔与 CBT--OHU 重写
机制，如何自然对应到这两个刻度。

\begin{remark}[广度支持集与崩溃强度]
前文“连续统崩溃$\Rightarrow$深度，同伦扭曲$\Rightarrow$广度”的备注中，
$\mathcal{B}_{\mathrm{CH}}(t)$、$B_{\mathrm{CH}}(t)$、$\mathcal{B}_{\mathrm{CH+op}}(t)$
沿用主稿 GZ-CBT 的诊断层口径，描述的是\emph{沿轨迹随时间变化的崩溃强度}（非负标量）。

而本节讨论“广度”时，我们需要一个\emph{静态的层级支持集}来记录：
对给定的 law-object~$L$，
JacGap 与同伦数据在\emph{哪些 CH 层级}上是“可见且活跃”的。
为避免与 $B_{\mathrm{CH}}(t)$ 混淆，本节将该支持集统一记为
$S_{\mathrm{CH}}(L)$（此前草稿中曾写作 $B_{\mathrm{CH}}(L)$，现已更名）。

两类量的关系是互补而非对立：连续统崩溃的强度（$\mathcal{B}_{\mathrm{CH}}$）
通常决定“修复要走多深”（步数/塔阶/扇区距离），而层级支持集 $S_{\mathrm{CH}}(L)$
与图表支持集 $B_{\mathrm{chart}}(L)$ 则刻画“修复需要动用多宽”（涉及多少层/图表/方向）。

更具体地说，同伦扭曲在主稿中主要通过\emph{激活高阶层 $l_{\ge3}$ 与同伦图表（homotopy chart）}
来被“看见”：扭曲越复杂，往往意味着需要在更多 layer/chart 上显式审计与调节 JacGap，
从而 $S_{\mathrm{CH}}(L)$ 与 $B_{\mathrm{chart}}(L)$ 更容易变大，这对应你在前文所说的“广度”。
与此同时，连续统崩溃的强度更直接地推高 CBT--OHU 修复链的步数/塔阶/扇区迁移距离，
因而表现为“深度”增加；但两者并非互斥，连续统崩溃也可能迫使我们进入同伦图表
从而间接变宽，同伦扭曲也可能需要多轮 OHU 才能压回阈值从而间接变深。
\end{remark}

\subsection{扭曲广度：CH-layering 与 PFB-GNLA 四图表的横向刻度}

先把主丛语言与主稿语言做一次显式对齐：在本笔记前文的主丛/联络表述中，
\textbf{广度}对应于\emph{在同一“主丛层级/塔阶”上}，通过扩展纤维或结构群的独立方向
来吸收/表征同伦扭曲（例如 $G\leadsto \widehat G$、$\mathfrak g\leadsto\widehat{\mathfrak g}$，
或把证书分量方向做直和扩张）。也就是说：\emph{不必引入更高阶的新主丛，只是在同一层里把“横向方向”变多}，
其语义在主稿里自然被 layer/chart 维度所编码。

在 G-Framework 主稿中，与 Jacobiator-gap 相关的“横向广度”主要通过
两类结构体现：
\begin{itemize}
  \item CH-layering 给出的层级 $l_n$（尤其是 $l_2$ 与 $l_{\ge3}$）；
  \item 四个 PFB-GNLA 图表（strict / homotopy /
        variable functional--operator / original law-level）。
\end{itemize}

\paragraph{(1) CH-layering 下的层级宽度}

主稿中通过 layer functor $\mathsf{Layer}$ 定义了 CH-layering 结构，
对每个连续统型 law-object $L_{\mathrm{cont}}$ 给出层级分解
\[
  \mathsf{Layer}_{\mathrm{CH}}
  :
  \mathsf{Law}_{\mathrm{Cont}}
  \longrightarrow
  \mathsf{Law}_{l_2}\times\mathsf{Law}_{l_{\ge3}},
  \qquad
  L_{\mathrm{cont}}
  \longmapsto
  \bigl(\Pi_{l_2}(L_{\mathrm{cont}}),\Pi_{l_{\ge3}}(L_{\mathrm{cont}})\bigr).
\]
其中：
\begin{itemize}
  \item $l_2$ 层捕捉纯外延/基数信息（如 $2^{\aleph_0}$ 的取值及其变更）；
  \item $l_{\ge3}$ 层聚合了所有高阶同伦/语义层，对 $\ell_3^H$、
        $H$-twist 等同伦信息敏感；
  \item 更一般的 $l_n$ 层则在更细粒度上拆分不同类型的语义/同伦约束。
\end{itemize}

从“扭曲广度”的角度看，可以对给定 law-object $L$ 定义一个
CH-layering 支持集：
\[
  S_{\mathrm{CH}}(L)
  =
  \bigl\{\,n \,\bigm|\,
    \text{在层级 } l_n \text{ 中 JacGap 或相关同伦数据是“可见且活跃”的}
  \bigr\}.
\]
集合 $S_{\mathrm{CH}}(L)$ 的基数与结构，刻画了在多少个 CH 层级上
需要同时考虑 Jacobiator-gap 的影响——这就是意义上的“层级广度”。
例如：
\begin{itemize}
  \item 若 $S_{\mathrm{CH}}(L)=\{2\}$，则 JacGap 只在外延层产生约束，
        高阶同伦层对其“盲视”；
  \item 若 $S_{\mathrm{CH}}(L)=\{2,\ge3\}$，则说明同一个 JacGap
        同时在外延与高阶同伦层起作用，需要跨层协调；
  \item 若进一步细分 $l_{\ge3}$ 为多个具体层级 $l_3,l_4,\dots$，
        则 $S_{\mathrm{CH}}(L)$ 的结构可以更加精细地记录
        “哪些具体语义层参与了扭曲”。
\end{itemize}

\paragraph{(2) PFB-GNLA 四图表下的实现广度}

主稿中 PFB-GNLA 的四个图表为：
\begin{itemize}
  \item strict chart：严格 Jacobi 的 PFB-GNLA，JacGap 在这里被“压平”
        （只允许非常受控的残余）；
  \item homotopy chart：显式引入 $H$-twist、$\ell_3^H$ 等，JacGap
        被提升为同伦结构的一部分；
  \item variable functional--operator chart：参数方向升级为算子型泛函，
        JacGap 与路径/算子同伦动态耦合；
  \item original law-level chart：在 PL-PI 元数学层面直接操作 law-space，
        JacGap 作为原始 law-level 缺陷出现。
\end{itemize}

对给定的法则系统 $L$，我们可以考虑在这四个图表中“实际工作”的子集：
\[
  B_{\mathrm{chart}}(L)
  =
  \bigl\{\text{chart}\in
    \{\mathrm{strict},\mathrm{homo},\mathrm{vf\text{-}op},\mathrm{orig}\}
    \mid
    \text{JacGap 在该 chart 中被显式调节或审计}
  \bigr\}.
\]
当 $B_{\mathrm{chart}}(L)$ 越大，说明在更多的实现层上需要同时追踪
和控制 JacGap，这在本笔记的语言中对应“实现层面的扭曲广度”。

综上，可以把
\[
  B(L)
  =
  \bigl(S_{\mathrm{CH}}(L),B_{\mathrm{chart}}(L)\bigr)
\]
视为 G-Framework 主稿内部对“扭曲广度”的一个自然刻度：它同时记录
了“在多少个 CH 层级”和“在多少个 PFB-GNLA 实现图表”上，JacGap
是活跃且需要被处理的。它与本笔记前文“广度 = 涉及多少 law-level 方向/子结构”
在\emph{粗粒度口径}下对齐：主稿用 layer/chart 编码“哪些方向被动员”，
而不直接枚举证书分量集合 $J_\ast$。
因此 $B(L)$ 可以看作对 $J_\ast$ 的一个投影性摘要（保留 layer/chart 支持，遗忘分量细节）。

\subsection{扭曲深度：重写链、law-bundle 塔与 law-space 扇区迁移}

同样地，在本笔记的主丛/联络表述中，
\textbf{深度}对应于\emph{因连续统崩溃（CH breakdown）等“纵向故障”而被迫上升主丛层级}：
需要引入新的主丛包裹（或把 bundling 结构向更高阶展开），从而形成“塔”的推进
（例如从 $\mathcal{E}_n$ 提升到 $\mathcal{E}_{n+1}$，或在同伦塔上把 $\alpha_\ast$ 推高）。
换句话说：\emph{深度是“新主丛/新层级的引入与展开”}，它在主稿里对应为重写链长、law-bundle
塔阶数与扇区迁移距离等纵向刻度。

与“横向”的层级广度相对，主稿中关于“纵向”的扭曲深度主要体现在
以下三个互相关联的结构：
\begin{itemize}
  \item CBT--OHU 重写序列的步数与 JacGap-minimal 极限；
  \item law-bundle 塔 $\pi_n:\mathcal{E}_n\to B$ 的阶数；
  \item law-space 在 CH-energy/JacGap 约束下划分出的“区域/扇区”
        以及其间的迁移路径。
\end{itemize}

\paragraph{(1) CBT--OHU 重写序列中的链长刻度}

在处理 CH 与 JacGap 的联合“故障”时，主稿引入了 CBT--OHU 重写方案，
产生一族离散序列
\[
  \bigl(L_{\mathrm{cont}}^{(n)},L_{\mathrm{op}}^{(n)}\bigr)_{n\ge0},
\]
其中：
\begin{itemize}
  \item CBT 步在 law-space 中沿着 CH-energy 梯度及约束
        $\mathcal{C}_{\mathrm{adm}}$ 做“小步移动”，修正 CH 侧的
        断裂与曲率；
  \item OHU 步在 operator 侧完成 Jacobiator-gap 的同伦补全。
\end{itemize}

在理想情形下，该序列收敛到某个 JacGap 最小的极限
$L_{\mathrm{cont}}^{(\infty)}$。对给定容忍阈值
$(\theta_{\mathrm{CH}},\theta_{\mathrm{Jac}})$，可以定义
“达到可接受状态所需的最小步数”
\[
  n_\ast
  =
  \min\bigl\{\,
    n\mid
    \mathcal{B}_{\mathrm{CH+op}}^{(n)}
    \le\theta_{\mathrm{CH}},
    \ \mathsf{JacGap}\bigl(L_{\mathrm{cont}}^{(n)}\bigr)
    \le\theta_{\mathrm{Jac}}
  \bigr\},
\]
把 $n_\ast$ 理解为在 CBT--OHU 链上“向下走了多深”才把 JacGap 控制到
可接受范围内——这正是“扭曲深度”的一个离散刻度版本。

\paragraph{(2) law-bundle 塔阶数作为几何深度}

主稿中还引入了 law-bundle 塔
\[
  \pi_n:\mathcal{E}_n\to B,
\]
其中 $\mathcal{E}_n$ 是带有至多 $n$ 阶 law-connection 与 CH-layering
数据的 law-bundle。随着 $n$ 的增加，塔的每一层都编码了更多关于
law-space 几何与同伦的信息。

在 JacGap 语境下，可以把达到某一目标正则性所需的最小塔阶数
\[
  n_{\mathrm{tower}}
  =
  \min\bigl\{\,n\mid
    \text{在 }\mathcal{E}_n\text{ 中存在 JacGap-minimal 的提升}
  \bigr\}
\]
视为另一种“扭曲深度”刻度：只有在足够高阶的 law-bundle 层级上，
才能几何化地吸收或严格控制 Jacobiator-gap。

\paragraph{(3) law-space 扇区及其迁移深度}

在更全局的视角下，主稿通过 CH-energy 与 JacGap 等诊断量，把
law-space $\GZLS$ 划分为若干“扇区”（sector），例如：
\begin{itemize}
  \item CH-regular 且 JacGap 较小的“良性扇区”；
  \item CH breakdown 显著或 JacGap bounded away from $0$ 的“问题扇区”；
  \item CBT/OHU 需要重点介入的“临界扇区”等。
\end{itemize}
在 CBT--OHU 重写过程中，系统沿着 law-space 中的路径在这些扇区之间
迁移。若用一个简化的图模型（扇区为节点，扇区间允许的重写/同伦为
边），则可以定义
\[
  d_{\mathrm{sector}}(L_{\mathrm{init}},L_{\mathrm{target}})
\]
为从初始 law-object $L_{\mathrm{init}}$ 到某个 JacGap-minimal
$L_{\mathrm{target}}$ 的最短扇区路径长度（边数）。这一数值同样可以
作为“扭曲深度”的几何/图论刻度：它度量了需要跨越多少个“故障等级”
或“结构相态”，才能完成 JacGap 的修复或绕行。

\paragraph{(4) 与本笔记中 $\alpha_\ast$ 的对位}

把上述三个刻度 $(n_\ast,n_{\mathrm{tower}},d_{\mathrm{sector}})$
放在一起，可以把它们理解为主稿内部对“扭曲深度”的多种投影：
\begin{itemize}
  \item $n_\ast$ 反映的是算法/过程层面，在 CBT--OHU 链上走了多长；
  \item $n_{\mathrm{tower}}$ 反映的是几何/结构层面，需要提升到多高阶
        的 law-bundle 才能吸收 JacGap；
  \item $d_{\mathrm{sector}}$ 反映的是全局相图层面，需要跨越多少不同
        的 law-space 扇区才能到达“安全”区域。
\end{itemize}

在本笔记的极简模型中，我们用序数刻度 $\alpha_\ast$ 统一描述
“同伦扭曲深度”。从 G-Framework 的角度看，可以把 $\alpha_\ast$
视为上述三个深度刻度的某种综合或提升：当 CBT--OHU 重写、law-bundle
塔以及扇区迁移都被编码到主丛同伦塔
$\{L_\alpha\}_{\alpha<\kappa}$ 中时，$\alpha_\ast$ 恰好捕捉到了
“在主稿允许的所有层级与相态下，完成 JacGap 修复所需的最小扭曲层数”。
因此，尽管“深度”的主要驱动往往来自连续统侧的崩溃强度，它最终仍会在同伦塔坐标下
体现为 $\alpha_\ast$ 的增大（名称指的是刻度所在的表示空间，而非单一因果来源）。

\subsection{雅可比间隙作为指标与泛函簇}

在上述统一视角下，雅可比间隙 $\JacGap$ 具有双重含义：
\begin{itemize}
  \item \textbf{作为指标（index）：}
        它在 cohomology 层面上给出一个上同调类
        \[
          [\JacGap]\in H^{k+1}(X;\Lambda),
        \]
        表征泛拓扑下性质路径的不可收缩性，
        是一种“障碍大小”的量化指标；
  \item \textbf{作为泛函簇（family of functionals）：}
        在本笔记前文中，我们已经把
        $\JacGap(L)=(c_1(L),\dots,c_k(L))$ 理解为一组
        证书型不变量。每个分量 $c_i(L)$ 可以看成是
        作用在路径空间上的一个泛函
        \[
          c_i(L):\Gamma\longrightarrow\mathbb{R},
        \]
        它们共同构成了一个“泛函簇”，
        对路径的可行性、安全性、稳定性等维度给出评价。
\end{itemize}

因此，从 kernel\_reference 的泛拓扑/泛抽象代数视角回看，
雅可比间隙既是一个\emph{标记泛拓扑障碍的上同调指标}，
又是一个\emph{在路径空间上实施审计的泛函簇}。
这与本笔记早前建立的“路径偏好关系语义”等价性完全兼容：
Jacobiator-gap 既刻画了在 law-space 中“哪里不能走”，
又通过与 D 结构的对应告诉我们“在可行路径簇中该如何选择”。

\subsubsection*{标量指标与子结构级函数族的一致性}

上述分析可以进一步细化到“标量指标落地粒度”与“子结构级函数族”的
双重视角，并同时对 D 结构与雅可比间隙成立。

\begin{proposition}[D 结构与雅可比间隙的标量化与细化]
在本笔记的设定下，存在如下两类结构：
  \begin{enumerate}
    \item \textbf{标量指标落地：}
      D 结构与 Jacobiator-gap 都可以通过单调聚合得到系统级标量指标，
      用于粗粒度比较不同系统的“整体好坏”；
    \item \textbf{子结构级函数族：}
      D 结构与 Jacobiator-gap 都可以被分解为一族“子 D 结构”
      与“子雅可比间隙”，从而在更细粒度上控制和审计不同维度。
  \end{enumerate}
\end{proposition}

\begin{proof}[证明思路]
\textbf{(1) 标量指标落地：}
对固定系统 $S$，D 结构给出了路径评分函数
$f^{(D)}_S:\Gamma_S\to\mathbb{R}$，Jacobiator-gap 给出了
$f^{(\JacGap)}_S:\Gamma_S\to\mathbb{R}$，且二者在
$\mathsf{PathSem}$ 中通过严格单调函数重标度同构。
选取任意一个在实数上严格单调的聚合算子
\[
  \mathsf{Agg}:\mathbb{R}^{\Gamma_S}\longrightarrow\mathbb{R},
\]
例如
$\mathsf{Agg}(f)=\sup_{\gamma\in\Gamma_S} f(\gamma)$
或加权平均，都可以定义系统级标量指标
\[
  I_D(S)=\mathsf{Agg}\bigl(f^{(D)}_S\bigr),
  \qquad
  I_{\JacGap}(S)=\mathsf{Agg}\bigl(f^{(\JacGap)}_S\bigr).
\]
由于 $f^{(\JacGap)}_S=\phi\circ f^{(D)}_S$ 且 $\phi$、$\mathsf{Agg}$
均严格单调，$I_D(S)$ 与 $I_{\JacGap}(S)$ 在系统间比较时给出
完全一致的排序，因此 D 结构与 Jacobiator-gap 都可在系统层面
被视为\emph{标量指标}，用于粗粒度落地比较。

\textbf{(2) 子结构级函数族：}
在 D 结构一侧，逻辑性度量通常具有多维分解
\[
  \mathcal{L}_D(\gamma;w)
  =
  \sum_{i=1}^n w_i\,d_i(\gamma),
\]
可以把每一维 $(w_i,d_i)$ 看作一个\emph{子 D 结构} $D_i$，
其路径评分为
$f^{(D_i)}(\gamma)=w_i\,d_i(\gamma)$。
这样我们得到一个函数族
$\{f^{(D_i)}:\Gamma\to\mathbb{R}\}_{i=1}^n$，
它们的线性组合重构出原始的 $f^{(D)}$，并允许在每个维度上
独立进行权重调整与安全约束——这就是 D 结构在子结构级的细化控制。

在 Jacobiator-gap 一侧，$\JacGap(L)=(c_1(L),\dots,c_k(L))$
本身已经给出了自然的分量分解。每个分量 $c_j(L)$ 可以视作
与某一“子 law-level 结构” $L_j$ 或某个特定约束维度相关的
子雅可比间隙。对每个 $j$，我们可以定义子路径评分
\[
  f^{(\JacGap_j)}(\gamma)
  =
  F_j\bigl(c_j(L),\gamma\bigr),
\]
其中 $F_j$ 是在该维度上的局部证书函数。
函数族
$\{f^{(\JacGap_j)}:\Gamma\to\mathbb{R}\}_{j=1}^k$
共同决定了 $f^{(\JacGap)}$，并允许我们在每个分量上单独施加
“局部容忍度”或“局部优化目标”，从而实现对子雅可比间隙的精细控制。

结合前文关于路径语义范畴 $\mathsf{PathSem}$ 与系统范畴
$\mathsf{Sys}$ 的构造，可以把这些子 D 结构与子雅可比间隙理解为
从 $S$ 到若干“坐标系”对象的函子性投影：在对象层面对应
“选择某一维度/分量”，在态射层面对应“在该维度/分量上保持单调与
严格重标度相容”。因此，D 结构与 Jacobiator-gap 同时具备
标量化与细化成函数族的两种表现形式。
\end{proof}

\begin{remark}[函数族视角比向量更具解释力]
上述命题中的“子结构级函数族”视角在严格意义上比单纯的“向量视角”
保留了更多语义信息，从而更具解释力。形式化地说：
\begin{itemize}
  \item 设 $\mathsf{FunFam}$ 的对象为
        $(\Gamma,\{f_i:\Gamma\to\mathbb{R}\}_{i\in I})$，
        态射为保持每一分量严格单调重标度的映射；
  \item 设 $\mathsf{VecRep}$ 的对象为 $(\Gamma,f:\Gamma\to\mathbb{R})$，
        态射为保持严格单调重标度的映射；
  \item 定义遗忘函子
        \[
          U:\mathsf{FunFam}\longrightarrow\mathsf{VecRep},\quad
          U(\Gamma,\{f_i\})=(\Gamma,\sum_{i\in I} f_i).
        \]
\end{itemize}
则 $U$ 不是在对象等价类上的单射：存在两个非同构的函数族
  $(\Gamma,\{f_i\}_{i\in I})$ 与 $(\Gamma,\{g_j\}_{j\in J})$，
  它们在 $\mathsf{FunFam}$ 中不可同构（例如 $I\neq J$ 且
  无双射 $\sigma:I\to J$ 使 $f_i$ 与 $g_{\sigma(i)}$ 在
  $\mathsf{PathSem}$ 意义下等价），但满足
  $\sum_i f_i = \sum_j g_j$，从而在 $\mathsf{VecRep}$ 中
  被 $U$ 识别为同一个对象。也就是说，向量视角对应的是对
  函数族视角施加了一个“聚合商”，部分抹去了不同分量的语义区分能力。

因此，只要我们关心“哪一个维度/子结构贡献了哪一种语义效应”
这类问题，函数族视角在范畴意义上就是严格更细的表示，
也就比向量视角更具解释力。
\end{remark}

\section{双向构造的合成与严格等价性}

\subsection{自然等价类下的互逆性}

上面两个方向的构造表明：
\begin{itemize}
  \item 从 D 结构 $(E,H,w,\mathcal{L}_D)$ 出发，可以构造一个
        law-level 系统 $L_w$，使得其 Jacobiator-gap 语义与 D 结构语义
        等价；
  \item 从 Jacobiator-gap 证书体系 $(L,\JacGap)$ 出发，可以构造一个
        D 结构 $(E,H,w,\mathcal{L}_D)$，使得其逻辑性度量语义与
        Jacobiator-gap 语义等价。
\end{itemize}

注意到两个方向都只在评分函数之间允许\emph{严格单调的重标度变换}，
因此我们天然地在 D 结构与 Jacobiator-gap 体系上引入如下等价关系：
\begin{itemize}
  \item 若两个 D 结构在路径偏好意义下产生同一个偏好关系
        $\succ_D$，则视为语义等价；
  \item 若两个 Jacobiator-gap 体系在路径偏好意义下产生同一个
        偏好关系 $\succ_{\JacGap}$，则视为语义等价。
\end{itemize}

在这两个等价类空间上，上述双向构造互为逆映射，从而我们得到：

\begin{theorem}[雅可比间隙与 D 结构的语义等价性]
在路径偏好意义的自然等价类上，O3 理论中的 D 结构
  与 G-Framework 中基于 Jacobiator-gap 的 law-level 证书体系，
  存在一一对应：
  \[
    [D \text{-结构}]
    \;\cong\;
    [\text{Jacobiator-gap 证书体系}].
  \]
  换言之，二者在“对路径的决策语义”这一层面上是严格等价的。
\end{theorem}

\subsection{面向工程与理论的含义}

这一等价性具有两层含义：
\begin{itemize}
  \item \textbf{理论层面：}
        它表明 O3 理论中的 D 结构并非与 G-Framework 的 law-space
        几何脱节，而是在路径语义上与 Jacobiator-gap 所捕获的
        同伦结构完全对齐；
  \item \textbf{工程层面：}
        它允许在实现层面自由切换两种表述：
        在需要可解释性与法规审计时使用 Jacobiator-gap 证书；
        在需要直接驱动决策与自反性更新时使用 D 结构与逻辑性度量。
\end{itemize}

因此，“雅可比间隙”可以被视为 D 结构在 law-space 几何中的
一个\emph{坐标化呈现}，而 D 结构则是 Jacobiator-gap 在
O3 决策引擎中的\emph{可执行形态}。

\section{后续工作方向（草案）}

本笔记仅给出了语义层面的等价性证明草案，若要在严格的范畴论或
高阶同伦代数框架下完全形式化，还需要进一步工作，包括但不限于：
\begin{itemize}
  \item 为 D 结构引入明确的范畴对象与态射，刻画自反性更新的
        2-范畴结构；
  \item 在 G-Framework 中给出 Jacobiator-gap 证书的完整构造，
        明确其目标空间与范畴性质；
  \item 证明上文双向构造在范畴意义下给出的是一个等价（equivalence），
        而不仅仅是在集合层面的一一对应；
  \item 结合具体的 LLM 决策例子，给出可计算的 Jacobiator-gap
        与 D 结构参数的显式对应。
\end{itemize}

这些工作将留待后续版本与主稿的后续章节中系统展开。

\clearpage
\renewcommand{\refname}{\texorpdfstring{参考文献}{References}}
\begin{thebibliography}{9}
\bibitem{O3DStructureLogicalMetric}
Gao Zheng（高政）.
\emph{D结构与逻辑性度量：O3理论的基准驱动引擎}.
kernel\_reference 条目目录：
\url{sub_projects/open_meta_mathematical_theory/src/}\\
\url{kernel_reference/}\\
条目编号为 \texttt{1751598516}.

\bibitem{GaoZheng2025GFramework}
Gao, Z. (2025).
\newblock Meta-Mathematical Theory based on Pan-Logic Analysis and
Pan-Iterative Analysis (GaoZheng G-Framework) and the
Principal-Bundle-Based Generalized Noncommutative Lie Algebra
(GaoZheng G-Algebra).
\newblock In \emph{Meta-Mathematical Theory based on Pan-Logic Analysis and
Pan-Iterative Analysis (GaoZheng G-Framework) and the
Principal-Bundle-Based Generalized Noncommutative Lie Algebra
(GaoZheng G-Algebra): An Integrated Construction (v1.0)}.
\newblock Zenodo.
\newblock \url{https://doi.org/10.5281/zenodo.17651584}.


\bibitem{KernelReferenceGenericTopology}
Gao Zheng（高政）.
\emph{面向数学结构的泛拓扑的性变态射与泛抽象代数的性变算子：基于D结构驱动的数学结构性质变化}.
kernel\_reference 条目目录：
\url{sub_projects/open_meta_mathematical_theory/src/}\\
\url{kernel_reference/}\\
条目编号为 \texttt{1737014305}.

\bibitem{KernelReferenceSexualOperators}
Gao Zheng（高政）.
\emph{性变算子与性变态射的严格定义与统一视角}.
kernel\_reference 条目目录：
\url{sub_projects/open_meta_mathematical_theory/src/}\\
\url{kernel_reference/}\\
条目编号为 \texttt{1737014313}.
\end{thebibliography}

\clearpage
\end{CJK*}
\end{document}
